\documentclass[11pt,letterpaper]{article}
\usepackage[T1]{fontenc}
\usepackage[utf8]{inputenc}
\usepackage{lmodern}
\usepackage[margin=0.82in,headheight=15pt]{geometry}
\usepackage{amsmath,amssymb,mathtools}
\usepackage{array,booktabs,longtable,graphicx}
\usepackage{xcolor,microtype,fancyhdr,enumitem,needspace}
\usepackage{xurl}
\usepackage[colorlinks=true,linkcolor=blue!45!black,
  citecolor=blue!45!black,urlcolor=blue!45!black]{hyperref}
\hypersetup{pdftitle={Targeting improvements in G(k)},
  pdfauthor={Waring investigation -- working manuscript}}
\setlength{\parindent}{0pt}
\setlength{\parskip}{0.45em}
\setlength{\emergencystretch}{2em}
\setlist{itemsep=0.2em,topsep=0.35em}
\renewcommand{\arraystretch}{1.13}
\numberwithin{equation}{section}
\pagestyle{fancy}
\fancyhf{}
\fancyhead[L]{\small Targeting improvements in $G(k)$}
\fancyhead[R]{\small Working manuscript}
\fancyfoot[C]{\thepage}
\renewcommand{\headrulewidth}{0.3pt}
\begin{document}
\begin{titlepage}
\thispagestyle{empty}
\vspace*{0.6in}
{\large\color{blue!45!black}WARING INVESTIGATION\par}
\vspace{0.3in}
{\huge\bfseries Targeting improvements in $G(k)$\par}
\vspace{0.22in}
{\Large A two-hour screen and finite certificates\par}
\vspace{0.18in}
{\large 22 November 2026\par}
\vspace{0.6in}
\fcolorbox{blue!35!black}{blue!3}{\parbox{0.91\textwidth}{
\vspace{0.12in}
\textbf{The leading targets are $k=8$ and $k=9$.}\par
The follow-up calculations yield proposed working proofs of
\[
                     G(8)\le38,\qquad G(9)\le46.
\]
The exact moment recipes and the complete representation arguments are
included. The explicit primary-source comparison bounds located are
$39$ and $47$, respectively.
\vspace{0.12in}
}}
\par
\vspace{0.45in}
The screen accounts for the numerical optimization already present in the
literature, measures the distance to strict integer thresholds, and compares
the transfer of the earlier small-degree gains. It also repairs 44 entries
of the earlier large-degree numerical table using published inequalities
alone.
\par
\vfill
\textbf{Status.} This is a research working manuscript with exact finite
certificates and independent internal source and analytic checks. It is
prepared for further mathematical review; the literature search does not
establish a priority claim. Classical unconditional $G(k)$ is used throughout.
\end{titlepage}
\pagenumbering{roman}
\tableofcontents
\clearpage
\pagenumbering{arabic}

% Source section: report/overview.tex
\section{Outcome and research priorities}
\label{sec:outcome}

The strongest targets for transferring the recent small-degree improvements
are the pair $k=8,9$.  Testing that shortlist produced finite exact moment
certificates and complete proposed representation arguments for
\[
                         \boxed{G(8)\le38,\qquad G(9)\le46.}
\]
The explicit primary-source comparison bounds located in the literature
search are $G(8)\le39$ and $G(9)\le47$ \cite{W16}.  The proposed results
are recorded here as a working manuscript: the analytic arguments, source
hypotheses, and finite arithmetic have been reviewed within this
investigation, and the complete recipes are supplied for independent review.
The literature search does not establish priority.

\begin{center}
\small
\begin{tabular}{lrr}
\toprule
 & Eighth powers & Ninth powers\\
\midrule
Explicit published comparison & $39$ & $47$\\
Proposed bound & $38$ & $46$\\
Ordinary / smooth variables & $8+30$ & $11+35$\\
Fixed policies per round & $4$ & $4$\\
Finite rounds & $12$ & $6$\\
Total moment updates & $48$ & $24$\\
Conservative minor-arc saving & $7/100000$ & $1/2000$\\
\bottomrule
\end{tabular}
\end{center}

The ninth-power argument is the better starting point for a focused
manuscript review. It needs only sampled quadratic second moments, even
degree-eighteen coefficient moments, and the refined coefficient sampling
estimate proved in Section~\ref{sec:proof}. The eighth-power argument has
a smaller final margin and additionally uses the distribution of quadratic
rational heights and an auxiliary second moment. Its tightest parameter
condition has positive margin $4.5\cdot10^{-7}$; this is sufficient for
the finite small-smoothness convention, but must be retained explicitly.

\subsection{The decisive inequalities}

Throughout, $\Delta_v^{(k)}$ denotes a selected admissible upper excess,
not the unknown optimal exponent. It means
\[
 \int_0^1|g_k(\alpha;P,R)|^v\,d\alpha
          \ll P^{v-k+\Delta_v^{(k)}+\varepsilon},
\]
where $g_k$ is the complete smooth Weyl sum and the smoothness parameter is
sufficiently small. The readable sufficient pairs established by the
finite calculations are
\begin{align*}
 k=8:&\quad \Delta_{32}^{(8)}\le0.1789,
                  &\Delta_{34}^{(8)}\le0.11721,\\
 k=9:&\quad \Delta_{38}^{(9)}\le0.1851,
                  &\Delta_{40}^{(9)}\le0.1355.
\end{align*}
They satisfy, respectively,
\[
 \Delta_{32}^{(8)}+7\Delta_{34}^{(8)}<1,
 \qquad
 5\Delta_{38}^{(9)}+74\Delta_{40}^{(9)}<11.
\]
These strict inequalities are paired with the ordinary minor-arc endpoints
of \cite[Lemma~3.1]{W16}. The major arcs are handled separately; a minor-arc
inequality alone is not treated as a proof of a bound for $G(k)$.

\subsection{What changed in the target-selection decision}

The earlier authors had already optimized their methods numerically.
That concern was correct, and it shaped the comparison. We first
distinguished explicit literature bounds from incomplete numerical
specializations of the published inequalities. Doing so corrected 44
distinct entries of the earlier $21\le k\le100$ table, all by one integer,
using published rules alone. Several apparently tiny large-degree gaps
therefore ceased to be reasons to prioritize those degrees.

For degrees $5$--$20$, the screen compared the actual strict threshold with
several local gain profiles measured in the earlier degree-$5$, degree-$6$,
and degree-$7$ calculations. It accounted for the loss through the ordinary
and efficient differencing steps. These models consistently selected $8$
and $9$ as a leading pair. Their relative order depended on the choice of
calibration coordinate; the computation was not interpreted as a success
probability. The finite follow-up then supplied stronger evidence than the
initial extrapolation.

Degree $13$ remains interesting because of its small raw threshold gap.
The direct transfer model makes it substantially harder than $8$ or $9$,
and iterative feedback is a limitation of that model. Limited six-round
floating probes at both $10$ and $13$ improved their target residuals
without crossing zero. The extra quadratic-height resources gave no
additional gain at the selected degree-$13$ sizes. These observations
are not impossibility results. A further investigation of either degree
would benefit from examining a wider set of paths and lower-order inputs,
informed by the feedback analysis below.

A matched continuation tested the most immediate omitted path changes:
terminal zero decrements and deeper repeated terminals. It found a small
additional degree-ten gain and no degree-thirteen gain beyond shared
feedback. Even the optimistic terminal diagnostic remained above zero in
both enlarged retained families. These bounded follow-ups preserve the
priority of the degree-eight/nine pair.

\subsection{Why numerical reoptimization alone is not the explanation}

For the chosen ninth-power policies, we repeated the same six finite rounds
with specific analytic inputs replaced by valid alternatives. The statistic
below is $5\Delta_{38}+74\Delta_{40}-11$; it must be negative.

\begin{center}
\small
\begin{tabular}{lr}
\toprule
Inputs used with the same four fixed policies & Final statistic\\
\midrule
Restricted height means and refined coefficient sampling & $-0.049522969165$\\
Global smooth mean in place of each restricted height mean & $+0.444585706459$\\
Cube alternative in place of refined coefficient sampling & $+0.463140600000$\\
Both substitutions & $+0.463140600000$\\
\bottomrule
\end{tabular}
\end{center}

This is a controlled check of the chosen finite recipe. It shows that the
crossing in that recipe depends on the height-sensitive analytic inputs.
It does not show that every other parameter search with the alternative
inputs would fail. The feature requiring review is the adapted coefficient
estimate and its combination with published restricted-height means inside
the finite moment recurrence. The restricted-mean theorem itself is a
published input, not a novelty claim of this investigation.

\subsection{Reading this report}

Sections~\ref{sec:cross-k-screen} and~\ref{sec:literature-baseline} explain
the screen and its corrected baseline. Section~\ref{sec:proof} gives the
ninth-power argument. Section~\ref{sec:g8-input-audit} supplies the
eighth-power additions and its distinct major-arc pruning. The parameter
appendices print the finite recipes, and the accompanying Python/JSON
files reproduce the arithmetic with rational numbers. The complete
research archive also retains the exploratory data and the separate source
and numerical audits.

\par\bigskip\Needspace{10\baselineskip}

% Source section: thresholds/cross_k_synthesis.tex
\section{Selecting degrees by the distance to a strict improvement}
\label{sec:cross-k-screen}

The numerical screen favours the pair $k=8,9$ for transferring the recent
coefficient and rational-height estimates. Degree $13$ has an especially
small raw threshold gap, but the same calculation predicts substantially
greater attenuation of a new terminal saving there. The order of $8$ and $9$
within the leading pair is sensitive to the extrapolation model; their
advantage over the other small degrees is much more stable. The subsequent
finite calculations for both degrees cross their targets. This section
explains how those candidates were selected; their proofs and precise status
are treated separately below.

\subsection{Published optimisation and the comparison baseline}

The earlier tables already incorporate substantial numerical optimisation.
Vaughan--Wooley explicitly describe computer optimisation and iterative
refinement in \cite[Section~8]{VW00}; their Tables in Sections~9--22 supply
the degree-$7$ through degree-$20$ moment inputs. Wooley's later bounds for
intermediate powers optimise a mixture of ordinary and smooth sums
\cite[Table~2 and Section~3]{W16}. Br\"udern--Wooley optimise a different
annular expression for degrees $14$--$20$ \cite[Section~8]{BW23}. Consequently,
the question is whether an additional analytic input can cross an integer
threshold after reoptimisation.

For each full moment order $v$, let $\Delta_v$ denote a selected admissible
upper excess for the degree $k$ under discussion, so that
\[
 \int_0^1|g_k(\alpha;P,R)|^v\,d\alpha
       \ll P^{v-k+\Delta_v+\epsilon}.
\]
It is an upper value supplied by an available estimate, not the unknown
optimal excess. Equalities used to name table entries have this convention.
The baseline below uses the published tables, together with the previously
certified improvements for degrees $5,6,7$ and the degree-$14$ coefficient
curve produced during the preceding investigation. It does not incorporate
the newly selected degree-$8$ and degree-$9$ candidates. Published decimal inputs are treated
as the printed upper bounds, and all threshold arithmetic is rational.

\subsection{Two complementary threshold tests}

For $N$ variables, allocate an integer $a$ to ordinary Weyl sums and $N-a$
to smooth sums. Set
\[
 K=k(k+1),\qquad \omega=\frac aK,\qquad
 v=\frac{N-a}{1-\omega}.
\]
The ordinary minor-arc estimate in \cite[Lemma~3.1]{W16} gives the excess
\[
 E_k(N,a)=(1-\omega)\Delta_v-\omega.               \tag{S1}
\]
The test requires $E_k(N,a)<0$. Smooth real moments are obtained from the
lower convex envelope of the available moment bounds, and the search ranges
over every integer ordinary allocation. The associated major-arc pruning
conditions are checked separately. In particular, when $a\ge k+2$, the
ordinary major-arc bound alone suffices after bounding the smooth factors
trivially. For smaller $a$, the ordinary/smooth H\"older pruning criterion
or the applicable smooth major-arc estimate is retained as follows.

The small-ordinary-count gate requires a strong mean only on the major arcs,
not an epsilon-free full-circle smooth moment. Put
$p=k+2$, $t_0=\max\{7,\lfloor k/2\rfloor\}$, and
$U=(N-a)/(1-a/p)$ for $0<a<p$. In
\cite[Theorem~1.3]{BW24} we may take $\omega=\omega'=\Omega=1$:
indeed $t_0\ge7$ and $k\ge5$ give
\[
 1<\frac{2t_0+4}{t_0+10},\qquad
 1<\frac{2k}{k+4},\qquad 1<k/2.
\]
Taking that theorem's cutoff $Q=1$, and treating its remaining innermost
arcs trivially, gives the restricted strong mean
\[
 \int_{\mathfrak M(P)}|g(\alpha)|^U\,d\alpha
       \ll P^{U-k}\qquad(U>2t_0+4).
\]
Combining it with \cite[(4.1)]{W16} by H\"older yields
\[
 \int_{\mathfrak M\setminus\mathfrak N}|f|^a|g|^{N-a}
 \ll P^{N-k}W^{(\varepsilon-1/k)a/p}=o(P^{N-k}).
\]
Here the ordinary major arcs lie in the standard union
$\mathfrak M(P)$, and $W=\log\log P$. Every first zero-excess order
$v_0$ in the degree-$5$--$20$ input tables is at least $2t_0+4$, so the
implemented stricter test $U>v_0$ is conservative. In particular, the
degree-five value $U=19.6>18$ and degree-seven value
$U=32.625>18$ are justified without removing an epsilon loss from their
full-circle zero-excess moments. The direct smooth-major fallback follows
from the same theorem with $Q=W$, since the logarithmic core contains the
standard union $\mathfrak M(W)$.

The second test is the Br\"udern--Wooley quantity
\[
 \tau(k)=\max_{w\in\mathbb N}\frac{k-2\Delta_{2w}}{4w^2},
 \qquad G_0(k)=\min_{v\ge2}\left(v+\frac{\Delta_v}{\tau(k)}\right).
                                                               \tag{S2}
\]
Here the maximisation for $\tau$ is over \emph{even integer} moments;
real interpolation is allowed for $v$. Under the standard major-arc and
local-solubility conditions, $G_0(k)<N$ is sufficient for $G(k)\le N$
\cite[Theorem~6.2]{BW23}. Those additional variable-count conditions are
well below all the targets in this table.

As a sensitivity measure, suppose hypothetically that the relevant tail
excesses are replaced by $(1-\rho)\Delta_v$, with the low moments defining
$\tau$ held fixed. The required value of $\rho$ is minimised again over
the ordinary allocation or the annular moment. This is a numerical measure
of the target, not an assertion that such a uniform improvement is available.

\begin{table}[htbp]
\centering\small
\begin{tabular}{rrrrlr}
\hline
$k$ & Published $U$ & Target $U-1$ & Continuous boundary & Better test
    & Required $\rho$ (\%)\\
\hline
5  &17 &16 &16.99937 &ordinary &47.0375\\
6  &24 &23 &23.49889 &ordinary &20.0818\\
7  &31 &30 &30.39565 &ordinary &13.0999\\
8  &39 &38 &38.35774 &ordinary &7.1264\\
9  &47 &46 &46.23881 &ordinary &4.0403\\
10 &55 &54 &54.40901 &ordinary &5.5743\\
11 &63 &62 &62.81783 &ordinary &9.8330\\
12 &72 &71 &71.45942 &ordinary &4.7980\\
13 &81 &80 &80.25510 &ordinary &2.3031\\
14 &89 &88 &88.48172 &annular &3.8594\\
15 &97 &96 &96.45173 &annular &3.1258\\
16 &105&104&104.42744&annular &2.9627\\
17 &113&112&112.47480&annular &2.8820\\
18 &121&120&120.54599&annular &3.2998\\
19 &129&128&128.69117&annular &3.6978\\
20 &137&136&136.84395&annular &4.4786\\
\hline
\end{tabular}
\caption{The degree-$5$--$20$ screen before inserting the new degree-$8$ and
degree-$9$ candidates. For the ordinary test, the continuous boundary optimises the total
variable count with integer ordinary count; for the annular test it is $G_0$.
The sensitivity column is separately reoptimised, so its best allocation need
not equal the allocation defining the continuous boundary. All inequalities
at a successful integer target must be strict.}
\label{tab:cross-k-baselines}
\end{table}

The degree-$14$ row includes a small existing-input correction. The earlier
degree-$14$ curve lowers $\Delta_{26}$ from $4.0399382$ to
$4.038694654926$, giving
\[
 \tau(14)=\frac{1480652672537}{169000000000000},\qquad
 G_0(14)=88.48172292042932\ldots.
\]
The unrefined table instead gives $88.48696658968404\ldots$; neither
crosses $88$. We also tested all H\"older interpolations of the ordinary
endpoint with the lower ordinary moments supplied by Hua's lemma and the
replacement in \cite[Theorem~10]{Bou16}. None of the 60,114 exact candidate
vertices improves (S1) at these one-step targets.

\subsection{What the achieved improvements support}

Calibration uses pure smooth moments of the same degree. The later
two-frequency linear--quintic moments and unproved auxiliary targets are
excluded. The following examples show why a single percentage is unsuitable.

\begin{table}[htbp]
\centering\small
\begin{tabular}{rrrrr}
\hline
$k$ & Full moment $v$ & Published $\Delta_v$ & Achieved $\Delta_v$
 & Relative decrease (\%)\\
\hline
5&14&0.272728900&0.270071744055&0.9743\\
5&16&0.077362700&0.074048069211&4.2845\\
6&20&0.203005500&0.195441875472&3.7258\\
6&22&0.086041200&0.079950274887&7.0791\\
7&22&0.424297300&0.423201151219&0.2583\\
7&24&0.293220800&0.283697712880&3.2478\\
7&26&0.192537400&0.176197101178&8.4868\\
7&28&0.113929700&0.093564057962&17.8756\\
7&30&0.052884800&0.015452339536&70.7811\\
\hline
\end{tabular}
\caption{Achieved pure-moment calibration data from the preceding finite
certificates. The last row includes endgame recycling and is not used as an
independent terminal-gain benchmark in the principal transfer model.}
\label{tab:achieved-calibration}
\end{table}

The $41.53\%$ reduction previously quoted for the seventh-power endgame
compared two different best ordinary allocations. It is neither a universal
moment improvement nor a transferable theorem. Instead, the screen matches
the local gain profile against either the old excess $\Delta_v$ or the
position $v/v_0$, where $v_0$ is the first available optimal even moment.
A separate $v/k$ alignment is included to expose sensitivity to that choice.

For a fixed efficient-differencing path, the ordinary Weyl factor contributes
$2^{-(k-j-2)}$ and the $j-1$ backward Cauchy steps contribute
$2^{-(j-1)}$. Their product is $2^{-(k-3)}$. Relative to degree $7$, a
first-step terminal gain thus has the factor $2^{7-k}$, before any
additional interpolation loss. Iterative feedback can amplify the eventual
moment improvement, so this factor is a sensitivity model, not a universal
upper bound on what the method can achieve.

\begin{table}[htbp]
\centering\small
\begin{tabular}{rrrrr}
\hline
$k$ & Match $\Delta$, no damping & Match $\Delta$, damped
 & Match $v/v_0$, damped & Match $v/k$, damped\\
\hline
8  &0.6226&1.2452&1.2585&1.3210\\
9  &0.3123&1.2494&1.2629&2.2401\\
10 &0.4516&3.6130&4.0076&10.4345\\
12 &0.2999&9.5984&9.6873&149.1114\\
13 &0.1379&8.8254&9.0662&461.5283\\
14 &0.2159&27.6358&26.9163&1252.7776\\
\hline
\end{tabular}
\caption{Required multiples of the achieved degree-$7$ direct gain profile
at moments $22$--$28$, excluding recycled moments $30,32$. A multiple below
one means that the corresponding hypothetical profile crosses the target.
These numbers describe assumptions, not probabilities. For the larger
degrees the active $v/k$ coordinate lies beyond the measured calibration
range, so that column is evidence of extrapolation weakness rather than a
reliable quantitative prediction.}
\label{tab:transfer-sensitivity}
\end{table}

Repeating the damped comparison with the achieved degree-$5$ and degree-$6$
profiles again selects degrees $8$ and $9$ first, although their absolute
predictions are less optimistic. Under excess matching, the two required
multipliers are approximately $(30.69,30.39)$ for degree-$5$ calibration
and $(6.92,6.86)$ for degree-$6$ calibration; their attenuation factors are
normalised to their respective calibration degrees. Thus the leading pair
is robust, while its internal ordering and the predicted gain size are not.

\subsection{Concrete targets and the finite follow-up}

The most useful sufficient inequalities are small enough to state directly:
\begin{align*}
 G(8)\le38 &: \quad \Delta_{32}+7\Delta_{34}<1,
                  &&\text{eight ordinary, thirty smooth};\\
 G(8)\le38 &: \quad 4\Delta_{30}+57\Delta_{32}<11,
                  &&\text{eleven ordinary, twenty-seven smooth};\\
 G(9)\le46 &: \quad 5\Delta_{38}+74\Delta_{40}<11,
                  &&\text{eleven ordinary, thirty-five smooth};\\
 G(13)\le80 &: \quad154\Delta_{68}+9\Delta_{70}<19,
                  &&\text{nineteen ordinary, sixty-one smooth}.
\end{align*}
For example, keeping the neighbouring published moment fixed, the first
degree-$8$ condition needs $\Delta_{34}<0.115840557142857\ldots$ instead
of $0.1294925$. The degree-$9$ condition needs
$\Delta_{40}<0.135798243243243\ldots$ instead of $0.1420569$.
The degree-$13$ condition needs
$\Delta_{68}<0.117469146103896\ldots$ instead of $0.1203776$.
The last raw gap is smaller, but its direct terminal-gain attenuation is
sixteen times greater than at degree $9$.

The selected degree-$9$ follow-up is now much more concrete than an
extrapolation. The canonical calculation uses four fixed policies, six finite
rounds and 24 moment updates, with only quadratic second moments, basic
coefficient moments and the refined reciprocal estimate. Its exact outputs
include
\[
 \Delta_{38}=\frac{185097318519}{10^{12}},\qquad
 \Delta_{40}=\frac{1693410547}{12500000000},
\]
and therefore
\[
 5\Delta_{38}+74\Delta_{40}-11
       =-\frac{9904593833}{200000000000}<0.
\]
The simpler rounded pair $(0.1851,0.1355)$ yields minor-arc saving
$19/36000>1/2000$. The certificate, source audits and complete classical
representation argument are the basis for the proposed $G(9)\le46$
bound; a negative floating optimisation alone would not establish it.

Degree $8$ required a second, broader numerical pass. Its canonical
construction now uses four fixed policies, twelve rounds and 48 updates.
Besides the preceding coefficient inputs, it uses the short quadratic-height
distribution estimate and one auxiliary $F_6$ second moment. The exact outputs
\[
 \Delta_{32}=\frac{178844726003}{10^{12}},\qquad
 \Delta_{34}=\frac{117131705619}{10^{12}}
\]
give
\[
 \Delta_{32}+7\Delta_{34}-1
       =-\frac{154166833}{125000000000}<0.
\]
The rounded pair $(0.1789,0.11721)$ supplies minor saving $7/100000$.
Major pruning follows already from the strong smooth $42$nd moment recorded
in \cite[Table~2 and (4.2)]{W16}, using ordinary moment $12$ and smooth
moment $90$. The eleven-ordinary alternative remains positive; the selected
eight-ordinary allocation is material. The source and finite arithmetic
audits for these additional inputs are given separately.

A limited six-round degree-$10$ floating comparison gives
$\Delta_{44}\approx0.18963221323273$ and
$\Delta_{46}\approx0.14697054764181$. Its sufficient criterion
$4\Delta_{44}+44\Delta_{46}-7$ decreases from $0.4132352$ to
$0.22523294917056$, with matched earlier-resource control
$0.2275307765956$, and remains positive. This pilot is included only as a
comparison of numerical progress, not as a new admissible-moment certificate
or evidence that no stronger degree-$10$ policy exists.

A targeted degree-$13$ probe updates full orders $62$--$70$ for six
floating rounds. It uses a degree-$26$ coefficient curve generated by the
published recurrence and all nine degree-$26$ auxiliary seeds in
\cite[Section~23]{VW00}, with coefficient
orders through $100$. The selected profiles use orders only through $68$.
The result is
\[
 \Delta_{68}\approx0.119715969771185,\qquad
 \Delta_{70}\approx0.100437062515795.
\]
Its statistic $154\Delta_{68}+9\Delta_{70}-19$ decreases from
$0.4479019$ to $0.340192907405$, removing about $24.05\%$ of the
initial deficit without crossing zero. Both matched resource families
give the same result. At their selected depth-nine sizes, the quadratic
product exponent exceeds $2.5$, so the newer quadratic-height estimate
coincides with the older one. There is also a structural explanation:
since $\phi_i\le1/k$,
\[
 h=k-2-kT\ge k-j-2.
\]
Thus at every depth $j\le k-4$ one has $h\ge2$ and
$\ell\le(h+2)/2\le h$. Throughout this height range the newer fourth-moment
resource $\min\{2,\max(\ell,2\ell-h)\}$ equals the older
$\min\{2,\ell,h\}$. A gain from this quadratic-height distribution input
therefore requires one of the two deepest depths, $j\ge k-3$.
The selected degree-thirteen depth is $j=9=k-4$.
Grouped coefficient sampling is a separate input, but these selected
coefficient resources are already
limited by their Fourier-support cap, and grouped sampling does not
remove that cap. This explains the lack of an extra gain at those fixed
sizes. It is not an obstruction for every possible parameter choice.

The degree-$10$ and degree-$13$ pilots are limited comparisons with
different retained pools and budgets; the latter disables auxiliary
interpolation. They reinforce the need to distinguish a raw threshold
gap from actual improvement through a selected proof. They do not settle
a universal ranking between these two fallback targets.

The reproducible small-degree data are in
\path{small_degree_thresholds_exact.json},
\path{calibration_ensemble.json}, and
\path{perturbation_scenarios_exact.json}. Their generating scripts keep
the published baseline, achieved inputs, and hypothetical transfer models
separate. The broader degree-$21$--$100$ screen is discussed in the following
section.

% Source section: thresholds/feedback_sensitivity.tex
\subsection{What finite moment feedback can change}
\label{sec:feedback-sensitivity}

The direct factor $2^{7-k}$ is a first-step comparison, not a prediction for
an already iterated moment table. To check whether iteration automatically
removes that attenuation, consider the unrounded analytic map for a finite
family of fixed rational size parameters,
height recipes and paths, and hold its degree-$2k$ coefficient moments fixed.
Write $q_i=1-T_i$ for a policy's smooth scale exponents. As a function of the
degree-$k$ full moment table, its auxiliary terminal is monotone and has
supremum-norm Lipschitz constant at most $q_j$. Indeed, the minor terminal's
coefficient of a smooth full exponent is $2c\ell/k\le cq_j\le q_j$,
where $0\le c\le1$ is the optional outer interpolation weight. The major
terminal has coefficient at most $2/k\le q_j$ because $kq_j/2>1$.

At a backward counted step, the Lipschitz constant is bounded by
\[
 \frac{q_i}{2}+\frac12\max(q_{i+1},L_{i+1})\le q_i.
\]
The first differencing step has the same bound $q_1=1-\phi_1$.
Taking minima over fixed policies, or taking the lower H\"older convex
envelope, preserves monotonicity and the corresponding Lipschitz bound.
With untouched orders anchored at their seed values, the complete fixed
family therefore has contraction factor at most
\[
 \rho=1-\min_{\rm policies}\phi_1<1.
\]
Here the appropriate map is the H\"older envelope of the minimum of the
published seeds and all fixed policy outputs. The implementation's operation
of retaining the previous bound is equivalent to that anchored map along
its monotone iteration from the seeds; a literal unrestricted map
$x\mapsto\min(x,T(x))$ would only have the immediate Lipschitz bound one.

The upward rounding in the actual finite certificate is treated separately:
the ceiling map on a decimal lattice is discontinuous and is not itself
claimed to be contractive. Each output differs from the unrounded map by
less than $10^{-12}$. The ordinary perturbation recurrence therefore bounds
the accumulated discrepancy after any finite number of these fixed-policy
steps by $10^{-12}/\min\phi_1$. This rounding qualification has no effect on
the validity of the upward-rounded moment certificates.

Consequently a uniform direct perturbation of size $\delta$ in the policy
outputs can change the associated fixed-family fixed point by at most
\[
 \frac{\delta}{\min\phi_1}.
\]
In particular, a change of at most $\delta_\sigma$ in the normalized saving for the
completed quadratic product, before taking the ordinary-differencing root, contributes at most
$2^{-(k-3)}\delta_\sigma/\min\phi_1$ after this fixed-family feedback.
If $\phi_1$ remains comparable to $1/k$, this allowance grows only linearly
with $k$ relative to the exponentially attenuated direct input. There is no
universal $O(k)$ assertion for policies permitting arbitrarily small
$\phi_1$, and the argument does not cover a simultaneous change of the
coefficient-degree moment table or a different analytic architecture.

For the actual frozen four-policy maps, exact two-sided rational differences
at scales $10^{-16}$ and $10^{-18}$, evaluated at the last finite certificate
iterate, give identical local Jacobians. A separate exact branch audit
confirms this stronger interpretation: it propagates affine gradients through
every changing H\"older-hull comparison, positive part and maximum. All
nonidentical branch comparisons have strictly positive rational slack;
zero-slack ties have identical affine gradients. The same formulas hold in
open supremum-norm boxes of radius at least $7.5\times10^{-7}$ for degree
eight and $2.0\times10^{-5}$ for degree nine, with the other input exponents
and the coefficient table held fixed. These Jacobians therefore describe
the active affine branches in a verified neighborhood of that state.
They are not a claim that the same branches
remain active throughout a future iteration. Their
spectral radii are $0.44206521$ for degree eight and $0.44745875$ for degree
nine. The local response norms $\|(I-J)^{-1}\|_\infty$ are approximately
$3.88182$ and $3.51233$, respectively. These are formal responses with those
local branches fixed, not global bounds on future amplification. The separate
global fixed-family bounds just proved are
$8.63039$ and $9.51633$. The final finite gains in the four target moments
are between $1.79$ and $3.72$ times the first-round gains for degree eight,
and between $1.80$ and $3.00$ times those gains for degree nine. These are
sensitivity measurements of the existing certificates, not further exponent
optimizations or an extra premise of the representation proofs.

Thus feedback is material and must be allowed for in target selection, but
it does not automatically erase the direct attenuation within these fixed
families. The earlier degree-seven calibration already includes its own
finite feedback, so applying another full amplification factor to that
profile would double-count it; a comparison would need the ratio of feedback
effects across degrees. This supports the prominence of degrees eight and
nine without excluding degree thirteen. In particular, the screen does not
settle a universal third-place ranking between degrees ten and thirteen.

\par\bigskip\Needspace{10\baselineskip}

% Source section: thresholds/ideal_terminal_headroom.tex
\subsection{An optimistic terminal test for degrees ten and thirteen}
\label{sec:ideal-terminal-headroom}

The two positive pilot deficits leave a more specific question than a
comparison of gain percentages: do the retained paths still have enough
terminal room if their coefficient estimates were ideal? We tested only the
previously retained pools: 37 paths updating full orders $38,40,\ldots,50$
for degree ten, and 23 paths updating orders $62,64,\ldots,70$ for degree
thirteen. No paths were added. These pools use only decrements one and two;
they omit the zero-decrement paths that mattered in the degree-eight result.
Neither retained pool contains an auxiliary outer interpolation.

There is an exact optimistic bound for the resource system. For an even
coefficient order $p$, the diagonal lower bound gives
$\mu_p^{(2k)}\ge p/2$, hence
$\Delta_p^{(2k)}\ge2k-p/2$. Each refined, cube or grouped coefficient
resource has its support cap
\[
 b_{i,p}\le(2k-\Delta_p^{(2k)})\phi_i\le p\phi_i/2;
\]
the basic order-two resource satisfies the same inequality. The quadratic
resources satisfy $z_p\le p/2$. Keeping the individual factor-load
constraints but discarding the total H\"older-mass constraint therefore gives
\[
 S\le\frac12+\frac12\sum_i\phi_i=\frac{1+T}{2}<1.
\]
At the top height $\ell=k(1-T)/2$, substitution of this optimistic saving
gives the relaxed minor-terminal exponent
\[
 \log_P A+\mu_{2r_j}^{(k)}(1-T)
       -2^{-(k-j-2)}\frac{1+T}{2}.
\]
The usual auxiliary major terminal, first diagonal and all counted backward
steps are retained. The oracle minimizes their maximum over the size
parameters for each fixed path. It drops additional coefficient restrictions,
so its output is an optimistic numerical envelope, not an admissible moment
estimate.

For each degree we iterated the same anchored map in two directions. The
upper run starts at the published table. The lower run replaces only the
updated coordinates by the necessary lower growth limits
$\max\{v/2,v-k\}$. All other seed coordinates remain fixed at their
published values, while H\"older convexification remains available at every
query. At each step the updated coordinates are clipped between those lower
limits and their published seed values. For any fixed size parameters the
oracle is monotone in every input exponent; taking a minimum over the same
feasible size domain and paths preserves monotonicity, as does the H\"older
envelope. Consequently this anchored construction has an increasing lower
iteration and a decreasing upper iteration. Using a routine that merely
retains the preceding upper bound would incorrectly freeze the lower run;
the diagnostic reconstructs the anchored map afresh instead.

The observed criterion intervals are as follows. A negative value is needed
to cross the corresponding record threshold.
\begin{center}
\begin{tabular}{rrrr}
\toprule
$k$ & iteration & lower criterion & upper criterion\\
\midrule
10&1&$-6.978445711$&$0.346594686$\\
10&10&$-0.335332055$&$0.188149251$\\
10&20&$0.184631245$&$0.186332161$\\
10&40&$0.186329223384$&$0.186329224995$\\
13&1&$-18.239797088$&$0.416364755$\\
13&10&$-0.838338615$&$0.325702234$\\
13&20&$0.320330805$&$0.324410464$\\
13&38&$0.324407851717$&$0.324407875402$\\
\bottomrule
\end{tabular}
\end{center}
The degree-ten criterion is $4\Delta_{44}+44\Delta_{46}-7$; the
degree-thirteen criterion is $154\Delta_{68}+9\Delta_{70}-19$.
The final whole-curve widths are below $6.1\times10^{-10}$, and the
observed lower/upper monotonicity and ordering violations are zero.
The corresponding remaining normalized minor deficits are approximately
$0.0033878041$ and $0.0017824609$.

These are floating LP diagnostics, not exact dual certificates or a
numerical impossibility theorem. Within the retained families, even the
idealized resource budget leaves both monitored criteria positive.
The result points to a concrete requirement for the next investigation:
expand the path family, enlarge the lower-order moment feedback system, or
change the analytic endgame or another part of the proof architecture.
Improving coefficient norms alone at these retained paths and anchored
other moments has very limited room. This does not exclude degrees ten or
thirteen, and does not settle their ranking under such changes.

\par\bigskip\Needspace{10\baselineskip}

% Source section: report/extended_screen.tex
\subsection{Matched checks of zero decrements and deeper terminals}
\label{sec:zero-extension}

The degree-eight success makes omitted zero decrements a concrete limitation
of the first degree-ten and degree-thirteen path pools. We therefore made a
bounded follow-up in both degrees. Each control and enlarged family started
from the same corresponding six-round floating curve and ran six additional
rounds with the same analytic resources. The permitted path edits were only
to replace the final decrement by zero and/or append one or two copies of
the terminal half-order, subject to $j\le k-2$. Earlier path positions were
unchanged.

Both comparisons used four height intervals, at most four alternating LP
steps per seed, $H_*\ge P$, and nonincreasing $\phi_i$. The last condition
is a search restriction, not a requirement of the underlying differencing
theorem. At degree ten the original engine's auxiliary options remained
available at the deepest new paths, but none survived the initial optimistic
screen. At degree thirteen auxiliary interpolation remained disabled.
Consequently these are matched comparisons within each degree, with the
stated difference of search scope between degrees.

For degree ten, 752 structural and auxiliary variants were examined and
five competitive additions were retained, enlarging the pool from 37 to 42
paths. For degree thirteen, 77 new variants were examined and 35 were
retained, enlarging the pool from 23 to 58. The final results are:
\begin{center}
\small
\begin{tabular}{rrrr}
\toprule
$k$ & common starting criterion & original-pool continuation
    & enlarged-pool continuation\\
\midrule
10 & $0.225232949170$ & $0.212661362045$ & $0.207688713163$\\
13 & $0.340192907405$ & $0.330078310753$ & $0.330078310753$\\
\bottomrule
\end{tabular}
\end{center}
The criteria are $4\Delta_{44}+44\Delta_{46}-7$ and
$154\Delta_{68}+9\Delta_{70}-19$, respectively. They must be negative;
their normalizing denominators are 55 and 182, so comparisons of raw
criterion sizes should be made within a row.

Four new terminal-zero paths won in the degree-ten calculation. The
extension-specific criterion gain was $0.004972648881$, about $2.34\%$ of
the matched control's remaining deficit. At degree thirteen every update
in the enlarged arm agreed with the control, and all winners remained at
depth nine. The final candidate values were
\[
\begin{array}{c|cc}
k&\text{first target excess}&\text{second target excess}\\\hline
10&\Delta_{44}\approx0.189451458901246&
   \Delta_{46}\approx0.146588247217238\\
13&\Delta_{68}\approx0.119655729340607&
   \Delta_{70}\approx0.100343999144414.
\end{array}
\]
A final-state rescan of the same finite variant sets found no initially
omitted variant becoming competitive under the optimistic oracle. This
addresses the possibility that continued feedback invalidated the initial
pruning of the added paths. It does not examine different early zero
placements or a larger path-generation rule.

We also repeated the two-sided ideal-terminal diagnostic on the enlarged
retained pools, with the same published/floor starts and anchored other
moments. Both arms now imposed the common $H_*\ge P$ and nonincreasing-size
restrictions. These restrictions left the original-pool numerical endpoints
unchanged. The enlarged-pool endpoints were
\begin{align*}
k=10:&\quad [0.182106476908956,\ 0.182106478520893],\\
k=13:&\quad [0.324407851717098,\ 0.324407875401860].
\end{align*}
They remain positive after 40 and 38 iterations, respectively. The added
degree-ten paths have a small amount of additional ideal terminal room;
the degree-thirteen ideal endpoint is unchanged. In these enlarged families,
the optimistic terminal allowance still does not make either monitored
criterion negative. The all-allocation check below uses the actual pilot
curves, not these idealized curves.

Every calculation in this subsection is a floating diagnostic. Neither the
fitted moments nor the ideal-budget brackets enter the proposed degree-eight
or degree-nine proofs. The evidence strengthens the choice of that leading
pair while leaving other zero placements, auxiliary interpolations,
lower-order feedback and different analytic architectures open.

\paragraph{Reoptimizing the final variable split.}
The preceding paired tables deliberately compare fixed criteria. After the
probes, we reevaluated every integer ordinary/smooth allocation, the pure
endpoints, the audited lower-ordinary-moment interpolations, and the annular
test on the resulting curves. The degree-ten winner changes to thirteen
ordinary and forty-one smooth variables. Its sufficient condition is
\[
                   73\Delta_{46}+24\Delta_{48}<13.
\]
On the enlarged floating profile, rounded upward to twelve decimals for this
reevaluation, the excess is
\[
 \frac{73\Delta_{46}+24\Delta_{48}-13}{110}
       \approx0.003525903204527>0.
\]
Degree thirteen retains the nineteen-ordinary allocation, with excess
approximately $0.001813617092412>0$. No alternative tested endgame crosses
either target. An optional deterministic published tail extension leaves
the winning allocations and all crossing decisions unchanged; none of the
reported conclusions depends on that extension.

The same evaluation checks the immediately lower integers for the exact
canonical curves. At $N=37$ in degree eight the best tested ordinary split
has six ordinary variables and positive excess
$0.029180148516083\ldots$; at $N=45$ in degree nine it has nine ordinary
variables and positive excess $0.021925559384$. The frozen moments thus
yield the stated one-integer improvements within these endgame families,
with no hidden second reduction. These comparisons do not assert optimality
among all possible representation arguments. The complete rational
reevaluation is in \path{postprobe_endgame_check.json}; rational evaluation
of the degree-ten/thirteen floating inputs does not certify those inputs.

\par\bigskip\Needspace{10\baselineskip}

% Source section: literature/literature_baseline_section.tex
\section{Published benchmarks, numerical optimization, and baseline repair}
\label{sec:literature-baseline}

Throughout this section, $G(k)$ denotes the classical, unconditional Waring
number over the rational integers: every sufficiently large natural number
must be representable as a sum of at most $G(k)$ positive integral $k$th
powers. This is distinct from the asymptotic-formula threshold
$\widetilde G(k)$, almost-all results, mixed-power problems, and number-field
analogues. In particular, the exact equality $G(4)=16$ precludes a reduction
to fifteen fourth powers; a later primary treatment of this sharp result
is \cite{DKW05}.

\subsection{The explicit published comparison bounds}

The explicit finite-degree upper bounds used for comparison are as follows.
They are taken from the main results of \cite{VW94,VW95,W16,BW23}.

\begin{center}
\begin{tabular}{rrl@{\qquad}rrl}
\hline
$k$ & bound for $G(k)$ & source & $k$ & bound for $G(k)$ & source\\
\hline
5 & 17 & \cite{VW95} & 13 & 81 & \cite{W16}\\
6 & 24 & \cite{VW94} & 14 & 89 & \cite{BW23}\\
7 & 31 & \cite{W16} & 15 & 97 & \cite{BW23}\\
8 & 39 & \cite{W16} & 16 & 105 & \cite{BW23}\\
9 & 47 & \cite{W16} & 17 & 113 & \cite{BW23}\\
10 & 55 & \cite{W16} & 18 & 121 & \cite{BW23}\\
11 & 63 & \cite{W16} & 19 & 129 & \cite{BW23}\\
12 & 72 & \cite{W16} & 20 & 137 & \cite{BW23}\\
\hline
\end{tabular}
\end{center}

For general degrees, \cite[Theorem~1.1]{BW23} supplies the convenient
displayed estimate
\[
 G(k)\le \left\lceil k(\log k+4.20032)\right\rceil.
\]
The finite-degree consequences of its proof can be stronger than this
rounded formula. Accordingly, three levels are recorded separately in the
data: an explicitly printed bound; our earlier numerical specialization of
published inequalities; and the improved numerical specialization obtained
in the present audit. A reduction relative to the second level does not,
by itself, establish a literature-priority claim.

\subsection{What the earlier papers had already optimized}

The suggestion that the previous authors had already optimized their
choices is supported directly by the sources. Vaughan and Wooley describe
computer optimization of their iterative processes in \cite[\S8]{VW00}.
The tables report permissible exponents together with the processes and
parameters producing them. Wooley's discussion before
\cite[Lemma~1]{W00QD} also addresses the difficulty of choosing fractional
moment parameters and the role of extensive computation. The optimization
of the annular quantities appears explicitly in \cite[\S\S7--8]{BW23}.

There is a useful qualification. The auxiliary larger-degree tables in
\cite[\S23]{VW00} were deliberately left less sharp than possible because
approximate inputs sufficed for the paper's principal tables. These cover
the seventeen degrees
\[
22,24,26,28,30,32,34,36,38,40,56,60,64,68,72,76,80.
\]
All 243 printed auxiliary rows were transcribed and checked. Injecting
every row into the earlier recurrence, rather than just the final row of
each table, left its final integer bounds unchanged. Additional published
processes nevertheless change the baseline substantially.

\subsection{The exact threshold and the numerical audit}

Write $\Delta_v$ for a permissible smooth excess, so that the full
$v$th moment is bounded by $P^{v-k+\Delta_v+\varepsilon}$. The quantities in
\cite[(5.1), (6.11), Theorem~6.2]{BW23} are
\[
 \tau(k)=\max_{w\in\mathbb N}
       \frac{k-2\Delta_{2w}}{4w^2},
 \qquad
 G_0(k)=\min_{v\ge2}\left(v+\frac{\Delta_v}{\tau(k)}\right).
\]
The resulting integer bound is
\[
 G(k)\le\max\{\lfloor G_0(k)\rfloor+1,L(k)\},
 \qquad
 L(k)=\begin{cases}4k,&k\text{ is a power of two},\\2k+3,&\text{otherwise}.
 \end{cases}
\]
Thus an existing bound $B$ decreases by one when the strict inequality
$G_0(k)<B-1$ and the local gate are both satisfied. The relevant gap is
$G_0(k)-(B-1)$. The maximum defining $\tau(k)$ uses even integer moments:
the source's Lemma~5.1 does not authorize an arbitrary real moment in that
maximum. Real orders remain available for the residual moment $v$.

We rebuilt the larger-degree baseline from the fourth-moment estimate,
the published auxiliary seeds with a conservative $10^{-6}$ increment,
and the recurrence displayed in \cite[proof of Theorem~2.1, (2.3)]{W93}.
The fractional-moment rule \cite[Lemma~1]{W00QD} and the auxiliary-moment
schemes \cite[Theorem~3.4 and \S4]{VW00} were then applied finitely many
times. Floating computations selected parameter policies; a separate
replay derived every adopted moment afresh with exact rational arithmetic
and rounded upper exponents upward to a $10^{-12}$ lattice.

The result is \emph{44 distinct corrected degrees} in $21\le k\le100$,
each improving the earlier session's integer bound by exactly one.
There are 38 corrections from the quasidiagonal closure and six additional,
disjoint corrections from the coupled auxiliary-moment closure. This count
is a count of degrees, rather than iterations, updates, or repeated
improvements of the same degree. Representative exact witnesses are:

\begin{center}
\begin{tabular}{rrrr}
\hline
$k$ & earlier bound & corrected bound & certified upper value of $G_0(k)$\\
\hline
21 & 147 & 146 & 145.923280248744\\
22 & 155 & 154 & 153.948991980478\\
34 & 257 & 256 & 255.969323186303\\
53 & 427 & 426 & 425.989994622735\\
58 & 473 & 472 & 471.999831515935\\
72 & 604 & 603 & 602.996606947951\\
74 & 623 & 622 & 621.983958260818\\
87 & 748 & 747 & 746.494797042825\\
\hline
\end{tabular}
\end{center}

The complete list of corrections and all eighty original/corrected rows
are in the accompanying CSV files. Each row names the exact certificate
and the published rule family. These are finite numerical consequences of
existing theory; no global optimum of all published processes is asserted.

This repair materially affects candidate selection. For example, after the
first quasidiagonal pass, degree 72 had $G_0(72)=603.00621064\ldots$, only
$0.00621\ldots$ above the threshold for replacing 604 by 603. Applying the
additional published schemes already crosses that threshold. The gap to
the next lower integer is then approximately $0.99661$. The analogous
effect occurs at degrees 22,34,53,58,74. Fractional closeness before this
audit is therefore insufficient evidence that a new analytic input is
likely to lower a properly updated bound.

For the smaller-degree comparison, the ordinary endpoint in
\cite[Lemma~3.1]{W16} has moment order $K=k(k+1)$ and saving one.
The lower ordinary moments of \cite[Theorem~10]{Bou16}, together with
Hua moments and interpolation with that endpoint, were checked separately.
The finite exact audit examined 60,114 candidate vertices for degrees 5--20
and did not improve the chosen endpoint at the one-variable target. This
is a check of these particular published alternatives, rather than a
claim that every possible ordinary input has been optimized.

\subsection{The scope of the priority and transcription checks}

The record search was carried out on 22 November 2026. It included
targeted searches for $G(8)\le38$, $G(9)\le46$, the corresponding numbers
of eighth and ninth powers, and 2025--2026 developments, as well as the
current author-maintained publication index. No primary source asserting
either improved integer was located. The best explicit primary-source
benchmarks located remain $G(8)\le39$ and $G(9)\le47$ from \cite{W16}.
Search absence does not establish priority and cannot exclude an
unindexed paper or an unpublished computation.

The embedded source tables for the two pilots were independently checked
against \cite{VW00}. For the eighth-power pilot, all eighteen degree 8
rows and twenty-eight degree 16 rows match \S\S10 and 18; the latter run
through full moment order 60. For the ninth-power pilot, all twenty
degree 9 rows and twenty-nine degree 18 rows match \S\S11 and 20; the latter
run through full order 62. The final-table pointwise exponents and the
starting orders of the $D_s$ process also match. These checks establish
accurate transcription of the published inputs. The analytic deductions
using them require the separate proof audit presented elsewhere in this
report.

\par\bigskip\Needspace{10\baselineskip}

% Source section: literature/next_threshold_ranking_21_100.tex
\subsection{The next integer thresholds after published-rule closure}
An earlier version of the corrected integer-bound table deliberately retained
older exact witnesses on unchanged-integer rows.  Ranking those witnesses
missed substantial published-rule gains, for example at degrees $25$, $31$,
$62$, $78$, and $95$.  The delivered table has now been synchronized with
the strongest available exact witnesses, choosing the smaller exact $H$
whenever two certificates have the same integer bound.
We therefore compared all saved quasidiagonal and coupled A/B runs, including
the unchanged-integer degree-$23$ coupled run, against the corrected target
$B_k-1$.  No new parameter search was performed.  The ten closest saved rows
were then replayed with exact rational arithmetic on their frozen policies.
Their ranks agree with the floating screen:
\begin{center}
\begin{tabular}{rrrrrrr}
$k$ & $B_k$ & next target & $H$ & $H-(B_k-1)$ & $u$ & $v$\\
\hline
38 & 292 & 291 & 291.055417950 & 0.055417950 & 74 & 254 \\
78 & 661 & 660 & 660.087065591 & 0.087065591 & 154 & 582 \\
95 & 825 & 824 & 824.096996413 & 0.096996413 & 188 & 730 \\
31 & 231 & 230 & 230.129285912 & 0.129285912 & 60 & 200 \\
25 & 180 & 179 & 179.136637734 & 0.136637734 & 48 & 156 \\
62 & 510 & 509 & 509.147555008 & 0.147555008 & 122 & 448 \\
70 & 585 & 584 & 584.157836494 & 0.157836494 & 138 & 516 \\
80 & 680 & 679 & 679.158583779 & 0.158583779 & 158 & 600 \\
88 & 757 & 756 & 756.160944774 & 0.160944774 & 174 & 670 \\
99 & 864 & 863 & 863.162255693 & 0.162255693 & 196 & 766 \\
\end{tabular}
\end{center}
Displayed decimals are rounded approximations; the companion JSON contains
the exact fractions for $H$, the gap, and both moment inputs in the witness
\[
 \tau=\frac{k-2\Delta_u}{u^2},\qquad H=v+\frac{\Delta_v}{\tau}.
\]
Every displayed target passes the applicable local-solubility gate, and every
$u$ is an even integer.  Thus the remaining obstacle in this specific
published-rule endgame is the strict inequality $H<B_k-1$, not a hidden
local gate.  The smallest gap among all saved post-closure runs is about
$0.055417947$ at $k=38$; none is below $0.05$.  This is a statement about
the finite computations inspected, not a lower bound for the optimal $G_0$
or an assertion that further published-rule optimization cannot cross a
threshold.  The coupled run at $k=23$ was also replayed exactly for 60 rounds,
with zero invalid policies.  It certifies the stronger gap
$0.224776286$, improving the earlier quasidiagonal gap $0.434792864$
at the same integer bound, but remains outside this top ten.

For comparison only, the labeled one-step attenuation factor $2^{7-k}$
is already $2^{-18}$ at the smallest degree in this list, $k=25$, and is
$2^{-31}$ at the closest-threshold degree, $k=38$.  These factors are not
probabilities or impossibility bounds, and do not model every possible
feedback or change of proof architecture.  Under that stated extrapolation,
the fractional threshold gaps alone do not justify moving these degrees
ahead of the much lower-degree targets supported by direct computations.

\par\bigskip\Needspace{10\baselineskip}

% Source section: report/proof_core.tex
\section{A finite working proof for ninth powers}
\label{sec:proof}

The screening led to a concrete candidate that is stronger than a heuristic
projection.  This section records the analytic argument used with the exact
parameter file.  The proposed conclusion is
\[
                         G(9)\le 46.
\]
The published comparison bound is $G(9)\le47$ \cite{W16}.  The present
argument and its finite arithmetic have been checked within this investigation;
this document is a working manuscript, and the literature search does not
establish a claim of priority for its constituent estimates.

\subsection{Notation and the two moments that suffice}

Write $e(z)=\exp(2\pi iz)$, and put
\[
 \mathcal A(X,R)=\{n\in\mathbb N:n\le X,\ p\text{ prime},\ p\mid n
                                      \Longrightarrow p\le R\},\qquad
 g_k(\alpha;X,R)=\sum_{x\in\mathcal A(X,R)}e(\alpha x^k).
\]
For a full moment order $v$, the notation $\Delta_v^{(k)}$ means an
admissible excess in
\begin{equation}
 \int_0^1|g_k(\alpha;X,R)|^v\,d\alpha
       \ll X^{v-k+\Delta_v^{(k)}+\varepsilon}.
 \label{eq:admissible}
\end{equation}
The assertion is uniform for a sufficiently small smoothness parameter,
depending on the finitely many orders and estimates being used.  We use
$\mu_v^{(k)}=v-k+\Delta_v^{(k)}$ for the full exponent.  This differs from
the half-order notation $\lambda_s$ in \cite{VW00}: $\mu_{2s}=\lambda_s$.

The final chosen certificate uses only sampled quadratic \emph{second}
moments.  It yields the rational permissible upper values
\begin{align}
 \Delta_{38}^{(9)}&=0.185097318519, &
 \Delta_{40}^{(9)}&=0.135472843760. \label{eq:chosen-pair}
\end{align}
Equality in this display specifies the admissible upper values supplied by
the computation, not the unknown optimal exponents.  In particular, the
slightly larger, convenient bounds
\begin{equation}
 \Delta_{38}^{(9)}\le\frac{1851}{10000},\qquad
 \Delta_{40}^{(9)}\le\frac{271}{2000}
 \label{eq:coarse-pair}
\end{equation}
are sufficient.  No conjectural moment estimate, limiting iteration, or
nonintegral coefficient moment is a premise of the certificate.

\subsection{Efficient differences and their quadratic factorisation}

We use the Cartesian sums and modified differences of \cite[Section 2]{VW00}.
For a fixed $j\in\{5,6\}$, choose $0<\phi_i<1/9$ and set
\[
 M_i=P^{\phi_i},\quad H_i\asymp PM_i^{-9},\quad
 T_i=\sum_{a\le i}\phi_a,\quad Q_i\asymp P^{1-T_i}.
\]
Fixed constants depending on the degree and fixed powers of $R$ are included
in $P^\varepsilon$ losses.  Write $T=T_j$, $Q=Q_j$,
$\widetilde H_j=\prod_{i\le j}H_i$,
$\widetilde M_j=\prod_{i\le j}M_i$, and $A=P\widetilde H_j\widetilde M_j$.
The coefficient variables run through the actual smooth blocks
$\mathcal C_i=\mathcal A(P,R)\cap(M_i,M_iR]$.

Let $F_j(\alpha)$ be the auxiliary sum of \cite[(2.1)]{VW00}.  Its
centred difference polynomial has, with $d=9-j$,
\[
 \Psi_j(z)=2^j9!\prod_i h_i
 \left\{\frac{z^d}{d!}+\frac{z^{d-2}}{6(d-2)!}
                       \sum_i h_i^2m_i^{18}
             +\text{terms of degree at most }d-4\right\}.
\]
Apply $r=d-2=7-j$ ordinary forward differences, with positive shifts
$u_1,\ldots,u_r$.  If $\xi=2z+\sum_lu_l$ and
$\Gamma=\prod_i h_i\prod_lu_l$, the resulting polynomial is exactly
\begin{equation}
 \frac{2^j9!}{24}\Gamma
       \left(3\xi^2+\sum_lu_l^2+4\sum_i h_i^2m_i^{18}\right).
 \label{eq:quadratic-factorisation}
\end{equation}
This is a difference of a cubic when $j=6$, and a second difference of a
quartic when $j=5$.  Terms of lower $z$ degree disappear.  Thus each
$m_i$ occurs in a separate degree-eighteen coefficient sum.

It is essential to perform these ordinary differences before separating
the coefficient factors.  Cauchy's inequality over the existing parameter
tuple gives the unseparated Weyl bound
\begin{equation}
 |F_j|^{2^r}\ll P^{2^r-1}B^{2^r}
       +P^{2^r-r-1}B^{2^r-1}|\mathcal G_r|,
 \qquad B=\widetilde H_j\widetilde M_j.
 \label{eq:unseparated}
\end{equation}
For completeness, if $\mathcal G_u$ contains $u$ ordinary shifts, Cauchy
over their tuple and the original $(\mathbf h,\mathbf m)$ gives
$|\mathcal G_u|^2\ll P^{2u+1}B^2+P^uB|\mathcal G_{u+1}|$.
Induction gives \eqref{eq:unseparated}.  It introduces no independent
copies of the original $m_i$ tuple.  The source instance
\cite[(12.17)]{VW95} illustrates this distinction.

The centered polynomial is the definition of $\Psi_j$; it is not a change
of variable in the defining sum.  In \cite[(2.1)]{VW00} the original range
$1\le z\le P_j$ is Cartesian and independent of both $\mathbf h$ and
$\mathbf m$.  After the ordinary differences the remaining range is
$1\le z\le P_j-\sum_lu_l$, as in the proof of \cite[Lemma 3.1]{VW00}.
Thus the interval for $\xi=2z+\sum_lu_l$ and its residue class
$\xi\equiv\sum_lu_l\pmod2$ depend only on the ordinary shifts, and are
independent of every coefficient variable $m_i$.  This independence
preserves the coefficient cancellation in the next display.

Fourier completion of this interval and residue class costs $O(\log P)$.
For a fixed real linear modulation $t$, it reduces the bound for
$\mathcal G_r$ to
\begin{equation}
 \sum_{\mathbf h,\mathbf u}
      |Z_t(c_0\Gamma\alpha)|
      \prod_{i=1}^j|L_i(c_i\Gamma h_i^2\alpha)|,
 \label{eq:completed-product}
\end{equation}
where $c_i$ are fixed nonzero integers,
\[
 Z_t(\beta)=\sum_{|z|\le CP}e(\beta z^2+tz),\qquad
 L_i(\beta)=\sum_{m\in\mathcal C_i}e(\beta m^{18}).
\]
The completion kernel has a common integrable majorant, independent of
the shift-dependent interval.  All estimates below are uniform in $t$,
so the common completion is legitimate.

Define
\[
 H_*\asymp P^{7-j}\widetilde H_j=P^{7-9T},\qquad
 V=P^2H_*\asymp Q^9.
\]
Every product of the fixed number of positive integer shifts has divisor
multiplicity $P^\varepsilon$.  Products may therefore be enlarged to the
full interval of length $H_*$ in nonnegative moment bounds.

\subsection{Intrinsic rational height and the quadratic second moment}

Put
\[
 \mathcal H_Q(\alpha)=\min_{\substack{q\ge1,\ a\in\mathbb Z\\(a,q)=1}}
                     \max(q,Q^9|q\alpha-a|).
\]
Its sublevel sets are precisely the standard major-arc unions, up to
fixed endpoint constants.  Dirichlet approximation gives
$\mathcal H_Q(\alpha)\ll Q^{9/2}$.  On the original auxiliary minor
arcs of \cite[Definition 4.9]{VW95},
$\mathcal H_Q(\alpha)\gg P^{1-O(\eta)}$ when $R=P^\eta$.
Indeed, a standard major-arc union of cutoff $cPR^{-C}$ lies inside the
original auxiliary major arcs.  Equality of the two arc systems is not
required.

We decompose the minor arcs into dyadic height shells
$\mathcal H_Q(\alpha)\asymp P^\ell$.  For exponent bookkeeping the
range is
\begin{equation}
                    1\le\ell\le \frac92(1-T).
 \label{eq:height-range}
\end{equation}
The actual lower endpoint differs by $O(\eta)$, which is absorbed in the
admissibility convention after the finite parameters are fixed.

For the four chosen policies, the exponents of $H_*$ are
\[
 2.00000221,\quad 1.88634850,\quad 1.84146280,\quad 2.44508497.
\]
In particular $H_*>P$.  Expansion of the quadratic second moment and the
divisor bound for $z^2-w^2$ give, uniformly in $t$,
\[
 \sum_{h\le H_*}|Z_t(c\alpha h)|^2
 \ll P^\varepsilon\left\{PH_*+
       \sum_{n\le C'P^2}\min(H_*,\|c\alpha n\|^{-1})\right\}.
\]
The extra zero differences $z=-w$ when the completed interval includes
negative integers are included in the $PH_*$ term. With a Dirichlet
approximant at conductor $V$, and $L=q+V|qc\alpha-a|$, the
reciprocal-sum estimate and transference in
\cite[Lemma 4.4, proof following (4.13)]{VW95} bound this by
\[
 P^\varepsilon\{PH_*+V/L+H_*+P^2+L\}\ll P^{1+\varepsilon}H_*.
\]
Here $L\gg P^{1-O(\eta)}$, $L\ll\sqrt V$, and $H_*\ge P$.
This proves the normalized quadratic saving $z_2=1$ needed by the
certificate, with no fourth-moment or small-$H_*$ extension.

\subsection{The refined coefficient sampling estimate}

We give the key deduction explicitly.  For coefficients supported in
$1\le n\le N$, let
\[
 E(\alpha)=\sum_{h\le H}\sum_{z\le J}
              \left|\sum_{n\le N}c(n)e(\alpha h^3zn)\right|^2,
 \qquad V_0=H^3JN.
\]
Choose a sufficiently small constant multiple of
\begin{equation}
               Y\le\min\{N,J,(JNH^{-3})^{1/2}\}.
 \label{eq:Y-range}
\end{equation}
The large-sieve estimate of \cite[Lemma 4.4]{VW95}, followed by the
initial steps of its Lemma 4.5, gives
\begin{equation}
 E(\alpha)\ll P^\varepsilon JN
       \left(\frac H Y+\sum_{h\text{ exceptional}}
                       \frac1{s+JN|\alpha h^3s-c|}\right)
                       \sum_n|c(n)|^2.
 \label{eq:exceptional}
\end{equation}
For each $h$, choose a reduced approximant $c/s$ with
$s\le JN/Y$ and $|\alpha h^3s-c|\le Y/(JN)$.
The exceptional condition is
$s+JN|\alpha h^3s-c|<Y/2$; the remaining reciprocal terms sum to
$O(H/Y)$.  The other large-sieve terms are also absorbed there by
\eqref{eq:Y-range}.

Choose independently a common reduced approximant $a/t$ with
$t\le\sqrt{V_0}$ and $|t\alpha-a|\le V_0^{-1/2}$.  For an
exceptional $h$,
\[
 |ah^3s-ct|
 <\frac{H^3Y}{2\sqrt{V_0}}+
      \frac{\sqrt{V_0}Y}{2JN}<1.
\]
The last inequality follows by taking the fixed constant in
\eqref{eq:Y-range} small enough.  Hence $ah^3s=ct$, so $s\mid t$
and $t_0=t/s$ divides $h^3$.  Let $b$ be the least positive integer
such that $t_0\mid b^3$.  Then $b\mid h$ and $t_0\le b^3$.
Writing $\beta=\alpha-a/t$, the contribution of this divisor class is
at most
\[
 \frac{t_0}{t}\sum_{n\le H/b}
                    (1+JN b^3n^3|\beta|)^{-1}.
\]
Nonempty classes have $b\le H$.  Trivial summation gives $H^3/t$;
comparison with $\sum_{n\ge1}n^{-3}$ gives
$O((tJN|\beta|)^{-1})$ when $\beta\ne0$.  Combining the bounds,
and treating $\beta=0$ by the first one, gives
\[
              \ll \frac{H^3}{t+V_0|t\alpha-a|}.
\]
There are at most $P^\varepsilon$ divisor classes.  Thus, with
$L=t+V_0|t\alpha-a|$,
\begin{equation}
 E(\alpha)\ll P^\varepsilon JHN
                  (Y^{-1}+H^2/L)\sum_n|c(n)|^2.
 \label{eq:refined-sampling}
\end{equation}
This deduction stops before the last step of the published proof that
identifies this internal rational approximant with an externally prescribed
one.  Its additional external matching conditions are consequently not
hypotheses of \eqref{eq:refined-sampling}.
The displayed conclusion of \cite[Lemma~4.5]{VW95} has the alternative
normalized term $L^{-1/3}$. The discrete class estimate above gives
$H^2/L$, which is smaller when $L>H^3$. This is an adaptation of that
published argument, and no priority claim is made for the variant.

Apply this estimate to the $p/2$ power of $L_i$, with $p$ a positive even
integer.  Grouping all shifts other than $h_i$ yields
\[
 J_i\asymp P^{7-j}\prod_{l\ne i}H_l,
 \quad J_iH_i^3M_i^{18}\asymp Q^9,
 \quad \log_PJ_i=6-9T+9\phi_i.
\]
The Fourier support is $N\asymp_p(M_iR)^{18}$.  Positivity of the
integer-power orthogonality count and the published degree-eighteen
moment give
\[
 \sum_n|c(n)|^2\ll P^\varepsilon M_i^{p-18+\Delta_p^{(18)}}.
\]
Each internal rational height in \eqref{eq:refined-sampling} is at
least a fixed multiple of the intrinsic height, with harmless powers
of $R$.  Fixed nonzero multipliers of $\alpha$ affect only constants:
reducing $a/(ct)$ gives $\mathcal H_Q(\alpha)\ll_c
t+Q^9|ct\alpha-a|$.  On the shell in \eqref{eq:height-range}, the
normalized coefficient saving is therefore
\begin{equation}
 b_{i,p}(\ell)=
 \min\left\{18\phi_i,6-9T+9\phi_i,
      \frac{3-9T+54\phi_i}{2},\ell-2+18\phi_i\right\}
           -\Delta_p^{(18)}\phi_i.
 \label{eq:b-refined}
\end{equation}
All preliminary size expressions in the selected policies are strictly
positive; the least exponent allowed for $Y$ is greater than $0.76$.

The published basic second-moment coefficient estimate is also used.
By \cite[Corollary 4.2.1 and Lemma 4.3]{VW95}, its normalized saving is
\begin{equation}
 b_{i,2}^{\rm basic}=
 \begin{cases}
 \phi_i,&j=5,\\[2pt]
 \min\{\phi_i,3-9T+36\phi_i,(1-T)/4,
                    (9(1-T)-35\phi_i)/6\},&j=6.
 \end{cases}
 \label{eq:b-basic}
\end{equation}
For the cited matching conditions take the Dirichlet cutoff
$X\asymp(Q^9/P)R^{9(9-j)}$. This retains the error cutoff in the
original auxiliary-arc definition. On the auxiliary minor arcs the
denominator exceeds the corresponding lower cutoff. The saving $b$ in
\eqref{eq:b-basic} is at most $\phi_i<1/9$, and the exact checks include
every additional size cap. In particular $9b<1$ has strict positive
margin, absorbing the displayed positive power of $R$ in the upper
matching condition $X\le Q^9P^{-9b}$ after smoothness is chosen small
enough. The smallest such exponent margin among the chosen resources is
$0.02982196$.

\subsection{H\"older allocation and the two auxiliary terminals}

On each of four fixed normalized height intervals choose a quadratic
second-moment weight $w_0\ge0$ and coefficient weights $w_{i,p}\ge0$,
where a coefficient resource may be either
\eqref{eq:b-refined} or \eqref{eq:b-basic}.  The certificate checks
\[
 2w_0\le1,\qquad \sum_p p w_{i,p}\le1\quad(i\le j),
 \qquad w_0+\sum_{i,p}w_{i,p}\le1.
\]
Unused factor loads are estimated trivially, and unused H\"older mass
is filled by the constant function one.  Consequently the saving in
\eqref{eq:completed-product} is
\[
                 S(\ell)=w_0+\sum_{i,p}w_{i,p}b_{i,p}(\ell).
\]
Insertion into \eqref{eq:unseparated} gives
\begin{equation}
 |F_j(\alpha)|\ll A P^{\varepsilon-
                         2^{j-7}\min(1,S(\ell))}.
 \label{eq:F-minor}
\end{equation}
The cap at one retains the ordinary-differencing diagonal contribution.

Br\"udern--Wooley's restricted major-arc mean estimate
\cite[Theorem 4.2]{BW23} states, for admissible $\Delta_v^{(9)}$ and
$1\le Z\le Q^{9/2}$,
\begin{equation}
 \int_{\mathcal H_Q(\alpha)\le Z}|g_9(\alpha;Q,R)|^v\,d\alpha
          \ll Q^{v-9+\varepsilon}Z^{2\Delta_v^{(9)}/9}.
 \label{eq:restricted-mean}
\end{equation}
At a fixed-factor endpoint the full mean supplies the same exponent.
The predecessor \cite[Lemma~2.3]{BW22PN} already supplies this restricted
mean in the moment range used by the present recipes. The estimate is a
published input; here it is applied inside the efficient-differencing
auxiliary recurrence.
Multiplying \eqref{eq:F-minor} by this shell mean gives the auxiliary
minor terminal, for remaining smooth order $v=2r_j$,
\begin{equation}
 A Q^{v-9}P^{\varepsilon+
  \max_\ell\{(2\Delta_v^{(9)}/9)\ell
                        -2^{j-7}\min(1,S(\ell))\}}.
 \label{eq:minor-terminal}
\end{equation}
Every coefficient saving is a minimum of affine functions of $\ell$,
hence concave.  For fixed nonnegative weights,
$\min(1,S(\ell))$ is concave.  The exponent in
\eqref{eq:minor-terminal} is convex on each assigned interval, so its
maximum is attained at an endpoint.  This justifies the finite endpoint
checks over the entire continuum of heights.

For the auxiliary major arcs, \cite[Lemmas 4.7--4.8, (4.25)--(4.26)]{VW95}
give the pointwise majorant
\[
 |F_j^*(\alpha)|\ll P^\varepsilon A
      \{q(1+Q^9|\alpha-a/q|)\}^{-1/(9-j)}.
\]
One can see the necessary divisor averaging directly by grouping
$\prod_i h_i$ and summing $b^{-1+1/(9-j)}$ over $b\mid q$.
Integrating its dyadic rational-height shells with
\eqref{eq:restricted-mean} gives the major terminal
\begin{equation}
 A Q^{v-9}P^{\varepsilon+
             \max(0,2\Delta_v^{(9)}/9-1/(9-j))}.
 \label{eq:major-terminal}
\end{equation}
The approximation error in Lemma 4.7 is at most
$A P^{-1/(9-j)+\varepsilon}$ because $q\le P$.
Its integral is also bounded by \eqref{eq:major-terminal}; the
approximation error is not discarded.  The required range $P\le Q^{9/2}$
holds with substantial strict slack for every policy.

\subsection{The finite counted recurrence}

All remaining smooth half-orders are integers.  Let
$r_1=s-1,r_2,\ldots,r_j$ denote a fixed path.  In exponent notation put
$q_i=1-T_i$ and $h_i=i-9T_i$.  Start $C$ at the maximum of the
major terminal exponent \eqref{eq:major-terminal} and all minor
endpoint exponents \eqref{eq:minor-terminal}.  For $i=j-1,\ldots,1$,
put $r=r_i$, $u=r_{i+1}$ and update
\begin{align}
 D_i&=1+h_i+T_{i+1}+\mu_{2u}^{(9)}q_{i+1},\\
 B_i&=\frac12\{\mu_{2(2r-u)}^{(9)}q_i+h_i+T_i
                                      +(2u-1)\phi_{i+1}\},\\
 C&\longleftarrow B_i+\frac12\max(D_i,C).
 \label{eq:counted-backward}
\end{align}
These are \cite[Lemma 2.2, including its corrected indices]{VW00}.
The condition $0<2r-u<2r$ is checked at every stage.  The first step
\cite[Lemma 2.1]{VW00} then supplies
\begin{equation}
 \mu_{2s}^{\rm new}=(2s-3)\phi_1+
  \max\{1+\phi_1+\mu_{2s-2}^{(9)}(1-\phi_1),C\}.
 \label{eq:first-step}
\end{equation}
The initial ordinary-square mixed mean dominates the pure even smooth
moment by inclusion of its nonnegative orthogonality count.  No
pointwise comparison between the ordinary and smooth sums is used.

The four paths are
\[
\begin{array}{ccl}
s&2s& (r_1,\ldots,r_j)\\\hline
17&34&(16,15,14,13,12,11)\\
18&36&(17,16,15,14,13,12)\\
19&38&(18,17,16,15,14,13)\\
20&40&(19,18,17,16,15).
\end{array}
\]
Their sizes and all interval weights are printed in the certificate
appendix and embedded in the standalone JSON parameter file.
We apply all four policies six times, always reading only exponents
available at the start of a round.  This is a finite induction with
24 updates.  Every output is rounded \emph{upward} to twelve decimal
places; H\"older convexification uses only already admissible values.
The seed moments are the published degree-nine and degree-eighteen
tables, checked digit for digit against \cite[Sections 11 and 20]{VW00}.
Only even coefficient orders up to 62 are used.  A single sufficiently
small $\eta$ accommodates every estimate and all endpoint perturbations.

The exact replay produces \eqref{eq:chosen-pair} and
\begin{equation}
 5\Delta_{38}^{(9)}+74\Delta_{40}^{(9)}-11
             =-\frac{9904593833}{200000000000}<0.
 \label{eq:exact-criterion}
\end{equation}
The numerical part is independently reproducible with the Python standard
library, without a numerical optimizer or the earlier workspaces.

\subsection{The classical representation argument}

Let $f(\alpha)=\sum_{x\le P}e(\alpha x^9)$ and
$g(\alpha)=g_9(\alpha;P,R)$.  Use the ordinary major/minor dissection
of \cite[Section 2]{W16}.  Its Lemma 3.1 includes the endpoint
$90=9(9+1)$ and gives
\[
                    \int_{\mathfrak m}|f|^{90}\ll P^{80+\varepsilon}.
\]
H\"older with weights $(11,5,74)/90$ on the moments
$(|f|^{90},|g|^{38},|g|^{40})$ is applicable because the weights sum
to one and $38(5/90)+40(74/90)=35$.  Thus
\[
 \int_{\mathfrak m}|f|^{11}|g|^{35}
 \ll P^{37+(5\Delta_{38}^{(9)}+74\Delta_{40}^{(9)}-11)/90+\varepsilon}
 \ll P^{37-19/36000+\varepsilon},
\]
where the final bound uses only \eqref{eq:coarse-pair}.
The saving $19/36000$ is greater than $1/2000$.

For $W=\log\log P$, let $\mathfrak N$ be the logarithmic major arcs
of \cite[Section 4]{W16}.  Its equation (4.1) includes $s=k+2=11$.
Consequently
\[
 \int_{\mathfrak M\setminus\mathfrak N}|f|^{11}|g|^{35}
 \le P^{35}\int_{\mathfrak M\setminus\mathfrak N}|f|^{11}
 \ll P^{37}W^{\varepsilon-1/9}=o(P^{37}).
\]
This step requires no additional optimal smooth moment.  The narrow
major-arc approximations in the same source give a positive main term
for $f^{11}g^{35}$.  The singular series is uniformly positive in the
cited sufficient range, since $46>4\cdot9$.  The singular integral
has its usual positive gamma-factor value.  Therefore the number of
representations of every sufficiently large integer as eleven ordinary
and thirty-five smooth ninth powers is positive.  This completes the
proposed proof of $G(9)\le46$ from the finite admissible-moment certificate.

\par\bigskip\Needspace{10\baselineskip}

% Source section: thresholds/g8_auxiliary_and_quadratic_audit.tex
\section{A finite working proof for eighth powers}
\label{sec:g8-input-audit}

The degree-eight calculation uses the coefficient argument of
Section~\ref{sec:proof} with degree $16$ in place of degree $18$.
Use the same Cartesian auxiliary sums, with $k=8$ and
$j\in\{4,5,6\}$, and apply $6-j$ ordinary differences to reach the
quadratic polynomial. Its exact factorisation is
\[
 \frac{2^j8!}{24}\Gamma
       \left(3\xi^2+\sum_lu_l^2+4\sum_i h_i^2m_i^{16}\right).
\]
The case $j=6$ needs no ordinary difference. The pointwise cap by one is
still a valid weakening in that case. The independence of the completed
quadratic interval from the $m_i$ is unchanged. The counted recurrence
\eqref{eq:counted-backward}--\eqref{eq:first-step} applies with every
degree-$9$ exponent and scale replaced by its degree-$8$ counterpart.
The terminal exponent is modified only when the auxiliary interpolation
below is used.

The four output orders and paths are
\[
\begin{array}{cl}
2s&(r_1,\ldots,r_j)\\\hline
28&(13,12,11,10,9,8)\\
30&(14,13,12,11,11)\\
32&(15,14,13,12,11)\\
34&(16,15,14,14).
\end{array}
\]
An inner zero decrement is permitted by the source condition
$0<2r_i-r_{i+1}<2r_i$; no strict decrease of the path is required.
The complete sizes and recipes are printed in
Appendix~\ref{app:g8-parameters}. Twelve beginning-of-round updates of
all four policies give a finite induction with 48 steps. Its final
admissible excesses include
\[
 \Delta_{32}^{(8)}=0.178844726003,\qquad
 \Delta_{34}^{(8)}=0.117131705619.
\]
As throughout, these equalities name permissible upper values rather than
unknown optimal exponents.

Two further analytic ingredients are
needed: a sampled quadratic fourth moment and a second moment of the
quadratic efficient-differencing auxiliary sum. This section supplies their
derivations and checks the ranges used by the canonical certificate
\texttt{g8\_exact\_certificate.json}. The arithmetic recurrence itself is
verified separately.

\subsection{Quadratic sums with arbitrary linear modulation}

For a conductor $V\ge1$, define the intrinsic maximum height
\[
 \mathcal H_V(\alpha)=
 \min_{(a,q)=1,\ q\ge1}\max\{q,V|q\alpha-a|\}.
\]
This is comparable, within a factor two, with the corresponding minimum of
$q+V|q\alpha-a|$. Dirichlet's theorem gives
$\mathcal H_V(\alpha)\ll\sqrt V$. A minimising maximum-height
approximant can be chosen with denominator-error product at most one, so the
arbitrary-approximant reciprocal-sum bound used in the proof of
\cite[Lemma~4.4]{VW95} applies. Equivalently, one may first use a Dirichlet
approximant and the explicit transference argument following (4.13) there.
Fixed nonzero integer multipliers of $\alpha$ change these heights only by
fixed factors, by reducing or lifting the corresponding rational fraction.

Let
\[
 Z_t(\beta)=\sum_{z\in I}e(c\beta z^2+tz),
\]
where $I$ is an integer interval of length $O(P)$ contained in
$[-CP,CP]$, $c\ne0$ is a fixed integer, and $t$ is any real number.
One ordinary difference removes the linear phase up to a factor of modulus
one. The standard reciprocal-sum estimate and its transference therefore give,
uniformly in $t$,
\[
 |Z_t(\beta)|^2\ll P^\epsilon
       \left(P+\frac{P^2}{B}+B\right),\qquad
 B=\mathcal H_{P^2}(\beta).
\]
Since $B\ll P$, this implies
\[
 |Z_t(\beta)|^4\ll P^\epsilon
                   \left(P^2+\frac{P^4}{B^2}\right).       \tag{Q1}
\]
In particular, no approximation of a quadratic Gauss sum that is uniform in
a growing leading coefficient is being assumed. The argument remains valid
when the real number $t$ depends on the sampled integer below.

\subsection{Distribution of short heights of multiples}

Put $V=HP^2$, $L=\mathcal H_V(\alpha)$ and, for $1\le Y\ll P$, let
\[
 N(Y)=\#\{1\le h\le H:
                 \mathcal H_{P^2}(h\alpha)\le Y\}.
\]
A witness to membership has $q\le Y$ and
$\|hq\alpha\|\le Y/P^2$. Grouping $n=hq$ costs a divisor factor
$P^\epsilon$, since all parameters have fixed polynomial size in $P$.
Set $A=HY$ and $J=P^2/Y$, so that $AJ=V$ independently of $Y$.
The indicator of $\|n\alpha\|\le1/J$ is at most
$J^{-1}\min\{J,\|n\alpha\|^{-1}\}$. Applying the reciprocal-sum
estimate with the same conductor-$V$ approximant for every $Y$ gives
\[
 N(Y)\ll P^\epsilon
       \left(\frac{HY^2}{P^2}+\frac{HY}{L}+1+\frac{LY}{P^2}\right).
                                                               \tag{Q2}
\]
Here the fractions in (Q2) are unaffected by replacing intrinsic maximum
height by sum-height, since those heights are comparable by fixed constants.

The short approximation lifts to the possibly unreduced fraction $a/(hq)$.
Reducing the fraction can only decrease both its denominator and error.
Consequently $L\ll HY$, so $N(Y)=0$ unless $Y\gg L/H$.
On the nonempty range the term $1$ in (Q2) is absorbed by $HY/L$.
Also $L^2\ll HP^2$ absorbs $LY/P^2$ into that same term. Thus
\[
 N(Y)\ll P^\epsilon\left(\frac{HY^2}{P^2}+\frac{HY}{L}\right).
                                                               \tag{Q3}
\]
Dyadic summation beginning at $Y\asymp\max\{1,L/H\}$ yields
\[
 \sum_{h\le H}\mathcal H_{P^2}(h\alpha)^{-2}
 \ll P^\epsilon\left(\frac H{P^2}+
                  \min\left\{\frac HL,\frac{H^2}{L^2}\right\}\right).
                                                               \tag{Q4}
\]
Combining (Q1) and (Q4), uniformly in arbitrary real modulations $t_h$, gives
\[
 \sum_{h\le H}|Z_{t_h}(h\alpha)|^4
 \ll P^\epsilon\left(HP^2+
             P^4\min\left\{\frac HL,\frac{H^2}{L^2}\right\}\right).
                                                               \tag{Q5}
\]
This proves the required fourth-moment estimate.

In the efficient-differencing application the sampled integer is a product
of finitely many positive shifts. Grouping their product costs another
divisor factor. Its range has
\[
 H_*=P^{k-j-2}\prod_{i=1}^jH_i=P^{k-2-kT},
 \qquad H_*P^2\asymp Q^k.
\]
Thus, if $H_*=P^h$ and $L\asymp P^\ell$, normalization by $H_*P^4$
gives the saving
\[
 z_4(\ell)=\min\{2,\max(\ell,2\ell-h)\}.            \tag{Q6}
\]
The selected policies do not insert this nonconcave expression into an
endpoint-only test. Instead they use the globally valid weaker resource
\[
 z_4^{\rm refined}(\ell)=\min\{2,2\ell-h\},           \tag{Q7}
\]
and occasionally the still weaker old resource
$\min\{2,\ell,h\}$. Both are bounded above by (Q6) and both are concave
as minima of affine functions. Their fourth-moment loads are included in the
same quadratic-factor H\"older constraint. Every frozen recipe is therefore
controlled on its whole height interval by the certified endpoint values.

The quadratic second-moment resource is also valid uniformly in modulation.
For the four actual policies, the values of $h=6-8T$ are
\[
 1.69428296,\quad1.83213176,\quad1.78352592,\quad2.39658688.
\]
They are all greater than one. Expansion of the second moment, the divisor
bound for differences of squares, and the reciprocal-sum estimate give
\[
 \sum_{h\le H_*}|Z_t(h\alpha)|^2
 \ll P^\epsilon\{PH_*+V/L+H_*+P^2+L\}
 \ll P^{\epsilon+O(\eta)}PH_*
\]
on the auxiliary minor region. The last step uses $L\gg P^{1-O(\eta)}$,
$L\ll\sqrt V$ and $H_*>P$. The normalized saving is one.
This second-moment expansion is uniform for a common real modulation $t$,
as supplied by completion of the quadratic interval. The fourth-moment
estimate (Q5) permits the stronger arbitrary-$t_h$ statement, which is not
needed for that completion.

\subsection{The auxiliary second moment at depth six}

For degree eight and $j=6$, the central polynomial has degree two.
Let $N_i$ count solutions of
\[
 \sum_{l=1}^i h_l^2(m_l^{16}-n_l^{16})=0
\]
in the coefficient ranges, and put
$D_i=H_iM_i$ and $U_i=H_iM_i^2$, suppressing powers of $R$ only as
the usual $P^\epsilon$ losses. For a fixed nonzero integer $b$, the equation
\[
 h^2(m^{16}-n^{16})=b
\]
has $O(P^\epsilon)$ solutions: choose $h^2\mid b$, choose the nonzero
divisor $m-n$ of $b/h^2$, and then use the strict monotonicity of the
remaining positive-variable difference polynomial. Splitting off $m_i=n_i$
therefore gives the one-block recurrence
\[
 N_i\ll P^\epsilon
       \left(D_iN_{i-1}+\prod_{l<i}U_l\right),\qquad
 N_1\ll P^\epsilon D_1.                                  \tag{A1}
\]
This is the elementary divisor argument in \cite[Lemma~3.8]{VW95}.
The off-diagonal term in (A1) is at most the desired diagonal product when
\[
 \prod_{l<i}M_l\le H_iM_i,
 \quad\text{equivalently}\quad T_i+(k-2)\phi_i\le1.       \tag{A2}
\]
It follows by finite induction that $N_6\ll P^\epsilon\prod D_i$.

Lemma~3.3 of \cite{VW95} explicitly includes the case $k-j=2$ and gives
\[
 \int_0^1|F_6(\alpha)|^2\,d\alpha
       \ll P^{1+\epsilon}N_6
       \ll P^{1+\epsilon}\prod_{i=1}^6H_iM_i.             \tag{A3}
\]
There is no odd-degree hypothesis in this quadratic case. In the canonical
full-order-$28$ policy, the exact margins in (A2), for $i=2,\ldots,6$, are
\[
 0.10069819,\quad0.02990766,\quad0.00000045,
       \quad0.03242985,\quad0.26812241.
\]
The smallest is positive but small; the smoothness parameter is chosen
sufficiently small after these fixed rational parameters. The resulting
exponent in (A3) is
\[
 1+\sum_i(1-7\phi_i)=3.23249759.
\]

The inner smooth order in this policy is $16$. With
\[
 a=\frac15,\qquad c=\frac45,\qquad\gamma=\frac35,
\]
one has $a+c=1$, $\gamma+2a=1$ and $20c=16$. Hence on each height shell
\[
 \int |F_6||g_Q|^{16}
 \le(\sup|F_6|)^\gamma
        \left(\int_0^1|F_6|^2\right)^a
        \left(\int_{\rm shell}|g_Q|^{20}\right)^c.         \tag{A4}
\]
The smooth factor in (A4) is estimated by \cite[Theorem~4.2]{BW23}; the
other factors use the pointwise terminal and (A3). This justifies exactly
the additional interpolation in the order-$28$ policy.

The auxiliary major contribution is retained separately with the original
smooth order $16$. Its published excess $1.8428621$ is smaller than
$k/(2(k-j))=2$, so the derived height-pruning bound has the optimal main
scale. Although \cite[Lemma~4.10]{VW95} restricts $j\le k-3$, that lemma is
not invoked here: Lemma~4.7 covers $j=k-2$, Lemma~4.8 covers quadratic
complete sums, and their pointwise averaged bound combined with
\cite[Theorem~4.2]{BW23} supplies the stated major estimate.

\subsection{Exact source-range checks for the remaining resources}

For degree-eight coefficient block $i$, put
\[
 \log_PJ_i=5-8T+8\phi_i,\quad
 \log_PH_i=1-8\phi_i,\quad\log_PN_i=16\phi_i.
\]
The conductor identity is $J_iH_i^3N_i\asymp Q^8$. The degree-sixteen
published moments in \cite[Section~18]{VW00} give the coefficient norm.
The refined and cube reciprocal resources follow, respectively, from
$H_i^3/L$ and $H_iL^{-1/3}$ in the divisor decomposition of
\cite[Lemma~4.5]{VW95}, with
\[
 Y\le\min\{N_i,J_i,(J_iN_iH_i^{-3})^{1/2}\}.
\]
Their normalized savings are the minima of the corresponding four affine
caps minus $\Delta^{(16)}_p\phi_i$. The code independently recomputes
these as physical exponent differences rather than importing the optimizer.

For the canonical twelve-round certificate, all 3600 coefficient-resource
endpoint occurrences pass. The minimum $Y$ exponent is $0.6936722$ and the
minimum $\phi_i$ is $0.03227716$. Every selected coefficient order is even
and covered by the published rows or the elementary low-order bounds below.
In addition to its printed rows starting
at full order six, the calculation uses the elementary full-order-four
exponent $2$; it follows by factoring $x^{16}-y^{16}$ and the divisor bound,
or from Hua's lemma. The order-zero and order-two exponents are $0$ and $1$.
The first coefficient moments retain the full
Lemma~4.3 minimum at depths five and six; the simpler Corollary~4.2.1
applies only at depth four. All H\"older masses and factor loads are
checked exactly. The audit also verifies all 96 occurrences of (Q7), all
48 occurrences of its older weaker form, and the 744 quadratic second-moment
occurrences. The finite epsilon/smoothness convention is chosen after the
fixed policies, preserving each displayed size condition.

This smoothness choice is a finite one. Each previously admissible moment is
uniform for $R\le X^{\eta_0}$ with a sufficiently small positive $\eta_0$.
When $X$ is replaced by a scale $Q_i$ or $M_iR$, its exponent relative to
$P$ is bounded away from zero in these fixed policies. Shrinking the common
$\eta$ therefore makes all seed and newly derived moment estimates valid
at every scale used in the forty-eight updates. The finitely many additional
powers of $R$ are absorbed by taking $\eta$ small relative to the required
$\epsilon$ and to the strict parameter margins, including $9/20000000$ in
(A2). The proof uses no limit of infinitely many iterations and does not
require an $\eta$ uniform over such a limit.

\subsection{The thirty-eight-variable endgame}

Assume the audited finite construction supplies admissible excesses satisfying
\[
 \Delta_{32}\le0.1789,\qquad\Delta_{34}\le0.11721.
\]
Take eight ordinary and thirty smooth variables. H\"older weights
$(1,1,7)/9$ on the moments $(|f|^{72},|g|^{32},|g|^{34})$ give
$f^8g^{30}$. The ordinary minor mean from \cite[Lemma~3.1]{W16} has exponent
$72-8-1=63$. Therefore the excess relative to $P^{30}$ is
\[
 \frac{\Delta_{32}+7\Delta_{34}-1}{9}
       \le-\frac7{100000}<0.                              \tag{E1}
\]

For major pruning no new smooth-major-arc result is needed. Table~2 of
\cite{W16} has $r=21$ for $k=8$, and its equation~(4.2) gives the strong
optimal mean
\[
 \int_0^1|g|^{42}\ll P^{34}
\]
without a $P^\epsilon$ loss. Consequently
$\int_0^1|g|^{90}\le P^{48}\int_0^1|g|^{42}\ll P^{82}$.
Using equation~(4.1) of \cite{W16} at ordinary moment $12$ and H\"older
weights $2/3,1/3$ yields
\[
 \int_{\mathfrak M\setminus\mathfrak N}|f|^8|g|^{30}
 \ll (P^4W^{\epsilon-1/8})^{2/3}(P^{82})^{1/3}
 \ll P^{30}W^{\epsilon-1/12}=o(P^{30}),
\]
where $W=\log\log P$. The source table and the absence of a power loss
in its strong $42$nd moment were checked in the original PDF.
The narrow-core approximation and the standard positive singular series
complete the usual representation argument, since $38\ge4k=32$.
Thus an admissible pair as in (E1) implies $G(8)\le38$.

The canonical arithmetic pair is
\[
 \Delta_{32}=\frac{178844726003}{10^{12}},\qquad
 \Delta_{34}=\frac{117131705619}{10^{12}},
\]
whose exact criterion and minor saving are
\[
 -\frac{154166833}{125000000000},\qquad
 \frac{154166833}{1125000000000},
\]
respectively. These figures are stronger than (E1).
The alternative eleven-ordinary allocation requires
$4\Delta_{30}+57\Delta_{32}<11$ and remains positive for this pair of
canonical profiles, so the eight-ordinary allocation is essential to the
present candidate.

\par\bigskip\Needspace{10\baselineskip}

% Source section: report/reproduction.tex
\section{Reproduction and the scope of the checks}
\label{sec:reproduction}

The two main certificates are independent of all earlier workspaces. Each
needs just its checker and its parameter JSON, together with Python~3's
standard library. Use the commands below with assertions enabled, as in a
normal Python invocation. The parameter files contain the published input
tables, fixed rational sizes, integer moment paths, and height-interval
H\"older weights. The checkers derive all output exponents afresh; they do
not accept optimized output exponents as premises.

\subsection{Ninth powers}
\begin{verbatim}
python exact_height_replay.py --fixed \
  --source k9_z2_only_fixed_parameters.json \
  --rounds 6 --out k9_z2_only_exact_certificate.json
\end{verbatim}
The final criterion is
\[
                   -9904593833/200000000000.
\]
The precise admissible excesses are $0.185097318519$ and $0.135472843760$.

\subsection{Eighth powers}
\begin{verbatim}
python g8_standalone_replay.py \
  --source g8_fixed_parameters.json \
  --rounds 12 --out g8_exact_certificate.json
\end{verbatim}
The final criterion is
\[
                    -154166833/125000000000.
\]
The precise admissible excesses are $0.178844726003$ and $0.117131705619$.

Both two-file runs were repeated in empty directories. Independent root
replays reproduced the full canonical output files byte for byte. The
physical-scale audits separately rebuilt the coefficient caps and checked
each frozen occurrence against the theorem hypotheses. These are distinct
from rerunning the optimizer.

A further implementation imports none of the canonical replay code.
It computes absolute sampled-moment exponents, forms the H\"older and
ordinary-Weyl products directly, and uses exhaustive two-point chords for
interpolation. It reproduced all 72 raw updates, all 960 height endpoints,
every backward recurrence stage, and both final criteria exactly. Its
standard-library source is \path{independent_physical_verifier.py}; the
complete comparison ledger is included in the package.

\subsection{What is certified, and what remains a matter of proof review}

Every adopted exponent is rounded upward to a twelve-decimal rational
lattice. Each round reads only the moment bounds available at its start.
The lower convex envelope supplies H\"older interpolation. All selected
coefficient moment orders are positive even integers drawn from the printed
tables or the stated elementary low-order bounds. Every height interval is covered, and endpoint
testing is justified by concavity of each selected saving and convexity of
the resulting exponent. In particular, the nonconcave envelope of the two
quadratic fourth-moment savings is never silently used in an endpoint test.

The two proof chains use finitely many estimates. One sufficiently small
positive smoothness exponent can therefore be chosen after all rational
parameters and all positive margins have been fixed. The source convention
allows the final $P^\varepsilon$ loss to be arbitrarily small. Choosing it
smaller than the displayed minor-arc savings preserves the strict bounds.

An arithmetic certificate verifies a finite numerical consequence of its
analytic rules. It cannot, by itself, establish those rules. The proofs
and source checks in this manuscript address that separate requirement;
an external mathematical review remains appropriate before describing
these working results as established new literature records. No claim of
priority for an individual auxiliary estimate is made.

\subsection{Other reproducible outputs}

The accompanying small-degree screen records exact thresholds, ordinary
allocations, annular witnesses, and calibration scenarios. Its prior
degree-$5$, degree-$6$, and degree-$14$ inputs are included locally, with
source hashes, so the screen no longer depends on earlier workspace paths.
The corrected $21$--$100$ table gives the old and current numerical
specializations and names the proof family supporting each correction.
Its 44 improving rows have separate exact certificates.

The exploratory degree-$10$ and degree-$13$ pilots and the fixed-policy substitution check
are retained as diagnostics. Their roles are limited: the pilots are
floating comparisons, and the substitution check fixes parameters rather than
reoptimizing the control cases. Neither is used as a premise of the two
proposed representation bounds.

The ideal-terminal diagnostic additionally records the two-sided floating
iterations in the retained degree-$10$ and degree-$13$ families, including
their narrow zero-decrement extensions. The next-threshold large-degree
shortlist also has fresh exact replays of its saved policies. The
fixed-policy feedback discussion has a different status: its contraction
bound is proved analytically, and its local affine branches are checked
with exact rational arithmetic in \path{feedback_branch_audit.py}. These
sensitivity statements concern the stated fixed families and do not bound
all possible methods for any degree.

\par\bigskip\Needspace{10\baselineskip}

\clearpage\appendix

% Source appendix: transfer/appendix_g9_parameters.tex
\section{Complete finite certificate for G(9)}
\label{app:g9-parameters}
This appendix prints the complete fixed recipe corresponding to the file
\par
\path{k9_z2_only_fixed_parameters.json}.
It uses four policies, 6 finite rounds, and 24 counted updates.
\subsection{Published input exponents}
The following decimals are the exact published upper exponents used as seeds.
The first column of each block is the half-order $s$; the listed value is
$\mu_{2s}^{(d)}$ in $\int_0^1|g_d|^{2s}\ll P^{\mu_{2s}^{(d)}+\varepsilon}$.
All three low-order values and all interpolation operations are independent
of the numerical selection runs: the additional seeds are
$\mu_0^{(d)}=0$, $\mu_2^{(d)}=1$, and $\mu_4^{(d)}=2$.
Only H\"older interpolation between known permissible exponents is used.
\begin{center}
\small
\begin{tabular}{rr@{\hspace{2.2cm}}rr}
\toprule
$s$ & $\mu_{2s}^{(9)}$ & $s$ & $\mu_{2s}^{(18)}$\\
\midrule
3 & 3.0358052 & 3 & 3.0050870\\
4 & 4.1822894 & 4 & 4.0484489\\
5 & 5.4197057 & 5 & 5.1469671\\
6 & 6.7383084 & 6 & 6.2893077\\
7 & 8.1356346 & 7 & 7.4676627\\
8 & 9.6039271 & 8 & 8.6816032\\
9 & 11.1425026 & 9 & 9.9332381\\
10 & 12.7463430 & 10 & 11.2185989\\
11 & 14.4105835 & 11 & 12.5388988\\
12 & 16.1292111 & 12 & 13.8936352\\
13 & 17.8959526 & 13 & 15.2823305\\
14 & 19.7055987 & 14 & 16.7049932\\
15 & 21.5507274 & 15 & 18.1609021\\
16 & 23.4269614 & 16 & 19.6497248\\
17 & 25.3292029 & 17 & 21.1706906\\
18 & 27.2520471 & 18 & 22.7231221\\
19 & 29.1901860 & 19 & 24.3060925\\
20 & 31.1420569 & 20 & 25.9186447\\
21 & 33.1033373 & 21 & 27.5596769\\
22 & 35.0727119 & 22 & 29.2280189\\
 &  & 23 & 30.9224199\\
 &  & 24 & 32.6415685\\
 &  & 25 & 34.3841221\\
 &  & 26 & 36.1487034\\
 &  & 27 & 37.9339369\\
 &  & 28 & 39.7384452\\
 &  & 29 & 41.5608774\\
 &  & 30 & 43.3999089\\
 &  & 31 & 45.2542644\\
\bottomrule
\end{tabular}
\end{center}
Source: Vaughan--Wooley, \emph{Further improvements in Waring's problem, IV}, Sections 11 and 20~\cite{VW00}.
The stored pointwise metadata are $\sigma_{9}=0.012183$ and $\sigma_{18}=0.005095$; these two values are not queried by the frozen recipes, whose moment orders remain inside the printed tables.
\subsection{The four exact policies}
The path is $(r_1,\ldots,r_j)$ in the counted recurrence.
Each size exponent below has eight decimal places and is exact.
Put $T=\sum_i\phi_i$, $h=k-2-kT$, and $\ell_{\max}=k(1-T)/2$.
A normalized interval $[a,b]$ in the weight tables means
$1+a(\ell_{\max}-1)\le\ell\le1+b(\ell_{\max}-1)$.
\paragraph{Full order $34$.}
$j=6$, $(r_1,\ldots,r_j)=(16,15,14,13,12,11)$.
\[ (\phi_1,\ldots,\phi_j)=\bigl(0.10508250,\ 0.10289747,\ 0.09976839,\ 0.09517079,\ 0.08734286,\ 0.06529330\bigr). \]
$T=0.55555531$, $h=2.00000221$.
The terminal uses the direct remaining smooth moment; there is no auxiliary interpolation.
\paragraph{Full order $36$.}
$j=6$, $(r_1,\ldots,r_j)=(17,16,15,14,13,12)$.
\[ (\phi_1,\ldots,\phi_j)=\bigl(0.10612499,\ 0.10413773,\ 0.10119188,\ 0.09645505,\ 0.08825004,\ 0.07202381\bigr). \]
$T=0.56818350$, $h=1.88634850$.
The terminal uses the direct remaining smooth moment; there is no auxiliary interpolation.
\paragraph{Full order $38$.}
$j=6$, $(r_1,\ldots,r_j)=(18,17,16,15,14,13)$.
\[ (\phi_1,\ldots,\phi_j)=\bigl(0.10702571,\ 0.10513564,\ 0.10230226,\ 0.09758589,\ 0.08900681,\ 0.07211449\bigr). \]
$T=0.57317080$, $h=1.84146280$.
The terminal uses the direct remaining smooth moment; there is no auxiliary interpolation.
\paragraph{Full order $40$.}
$j=5$, $(r_1,\ldots,r_j)=(19,18,17,16,15)$.
\[ (\phi_1,\ldots,\phi_j)=\bigl(0.10779756,\ 0.10604903,\ 0.10331138,\ 0.09871146,\ 0.09023224\bigr). \]
$T=0.50610167$, $h=2.44508497$.
The terminal uses the direct remaining smooth moment; there is no auxiliary interpolation.
\subsection{Complete rational H\"older allocations}
An entry $R_p[n]$ means the refined coefficient resource of full order $p$
with weight $n/10^6$.  The entry $B[n]$ is the published basic second-moment
resource with the same weight.  The column $L_i$ identifies the coefficient factor
(so its index is one greater than the zero-based JSON index).
$Z_2[n]$ denotes the sampled quadratic second-moment resource.
When two entries occur in one cell, both allocations are used and their
weights are added.  An empty cell has weight zero.  All weights are nonnegative.
The replay checks total H\"older mass at most one and total moment load at most one
for each factor.  These printed integer numerators, together with the preceding
exact size exponents, give the complete policy; no optimizer output is needed.
\begin{center}
\begin{minipage}{\linewidth}\centering
\footnotesize
\setlength{\tabcolsep}{3.3pt}
\renewcommand{\arraystretch}{1.22}
\textbf{Full order $34$: fixed weights}\par\smallskip
\begin{tabular}{cccccccc}
\toprule
$[a,b]$ & $Z$ & $L_1$ & $L_2$ & $L_3$ & $L_4$ & $L_5$ & $L_6$\\
\midrule
$[0,\frac{1}{4}]$ & $Z_2[499995]$ & $R_{44}[22727]$ & $R_{44}[22727]$ & $R_{46}[21738]$ & \shortstack[c]{$B[391976]$\\$R_{46}[4696]$} & $R_{50}[19999]$ & $R_{62}[16128]$\\
$[\frac{1}{4},\frac{1}{2}]$ & $Z_2[499995]$ & $R_{36}[27777]$ & $R_{36}[27777]$ & $R_{36}[27777]$ & \shortstack[c]{$B[362697]$\\$R_{38}[7226]$} & $R_{40}[24999]$ & $R_{46}[21738]$\\
$[\frac{1}{2},\frac{3}{4}]$ & $Z_2[499995]$ & $R_{26}[38461]$ & $R_{26}[38461]$ & $R_{26}[38461]$ & $R_{26}[38461]$ & $R_{28}[35713]$ & \shortstack[c]{$B[300958]$\\$R_{42}[9477]$}\\
$[\frac{3}{4},1]$ & $Z_2[499995]$ & $R_{14}[71427]$ & $R_{16}[62499]$ & $R_{18}[55554]$ & $R_{20}[49999]$ & $R_{26}[38461]$ & \shortstack[c]{$B[208154]$\\$R_{42}[13897]$}\\
\bottomrule
\end{tabular}
\end{minipage}
\end{center}
\medskip
\begin{center}
\begin{minipage}{\linewidth}\centering
\footnotesize
\setlength{\tabcolsep}{3.3pt}
\renewcommand{\arraystretch}{1.22}
\textbf{Full order $36$: fixed weights}\par\smallskip
\begin{tabular}{cccccccc}
\toprule
$[a,b]$ & $Z$ & $L_1$ & $L_2$ & $L_3$ & $L_4$ & $L_5$ & $L_6$\\
\midrule
$[0,\frac{1}{4}]$ & $Z_2[499995]$ & $R_{44}[22727]$ & $R_{44}[22727]$ & $R_{46}[21738]$ & $R_{46}[21738]$ & \shortstack[c]{$B[389397]$\\$R_{50}[4423]$} & $R_{58}[17241]$\\
$[\frac{1}{4},\frac{1}{2}]$ & $Z_2[499995]$ & $R_{36}[27777]$ & $R_{36}[27777]$ & $R_{36}[27777]$ & $R_{38}[26315]$ & \shortstack[c]{$B[361692]$\\$R_{40}[6915]$} & $R_{46}[21738]$\\
$[\frac{1}{2},\frac{3}{4}]$ & $Z_2[499995]$ & $R_{26}[38461]$ & $R_{28}[35713]$ & $R_{28}[35713]$ & $R_{28}[35713]$ & $R_{30}[33332]$ & \shortstack[c]{$B[311118]$\\$R_{38}[9940]$}\\
$[\frac{3}{4},1]$ & $Z_2[499995]$ & $R_{18}[55554]$ & $R_{18}[55554]$ & $R_{20}[49999]$ & $R_{24}[41666]$ & $R_{28}[35713]$ & \shortstack[c]{$B[248255]$\\$R_{38}[13249]$}\\
\bottomrule
\end{tabular}
\end{minipage}
\end{center}
\medskip
\begin{center}
\begin{minipage}{\linewidth}\centering
\footnotesize
\setlength{\tabcolsep}{3.3pt}
\renewcommand{\arraystretch}{1.22}
\textbf{Full order $38$: fixed weights}\par\smallskip
\begin{tabular}{cccccccc}
\toprule
$[a,b]$ & $Z$ & $L_1$ & $L_2$ & $L_3$ & $L_4$ & $L_5$ & $L_6$\\
\midrule
$[0,\frac{1}{4}]$ & $Z_2[499995]$ & $R_{44}[22727]$ & $R_{44}[22727]$ & $R_{46}[21738]$ & $R_{46}[21738]$ & \shortstack[c]{$B[389397]$\\$R_{50}[4423]$} & $R_{58}[17241]$\\
$[\frac{1}{4},\frac{1}{2}]$ & $Z_2[499995]$ & $R_{36}[27777]$ & $R_{36}[27777]$ & $R_{38}[26315]$ & $R_{38}[26315]$ & \shortstack[c]{$B[363231]$\\$R_{40}[6838]$} & $R_{46}[21738]$\\
$[\frac{1}{2},\frac{3}{4}]$ & $Z_2[499995]$ & $R_{28}[35713]$ & $R_{28}[35713]$ & $R_{28}[35713]$ & $R_{28}[35713]$ & $R_{30}[33332]$ & \shortstack[c]{$B[314533]$\\$R_{40}[9273]$}\\
$[\frac{3}{4},1]$ & $Z_2[499995]$ & $R_{18}[55554]$ & $R_{20}[49999]$ & $R_{22}[45454]$ & $R_{24}[41666]$ & $R_{28}[35713]$ & \shortstack[c]{$B[259585]$\\$R_{40}[12020]$}\\
\bottomrule
\end{tabular}
\end{minipage}
\end{center}
\medskip
\begin{center}
\begin{minipage}{\linewidth}\centering
\footnotesize
\setlength{\tabcolsep}{3.3pt}
\renewcommand{\arraystretch}{1.22}
\textbf{Full order $40$: fixed weights}\par\smallskip
\begin{tabular}{ccccccc}
\toprule
$[a,b]$ & $Z$ & $L_1$ & $L_2$ & $L_3$ & $L_4$ & $L_5$\\
\midrule
$[0,\frac{1}{4}]$ & $Z_2[499995]$ & \shortstack[c]{$B[403306]$\\$R_{42}[4604]$} & $R_{42}[23809]$ & $R_{42}[23809]$ & $R_{44}[22727]$ & $R_{46}[21738]$\\
$[\frac{1}{4},\frac{1}{2}]$ & $Z_2[499995]$ & \shortstack[c]{$B[365805]$\\$R_{30}[8945]$} & $R_{30}[33332]$ & $R_{32}[31249]$ & $R_{32}[31249]$ & $R_{34}[29411]$\\
$[\frac{1}{2},\frac{3}{4}]$ & $Z_2[499995]$ & \shortstack[c]{$B[187498]$\\$R_{14}[44642]$} & $R_{14}[71427]$ & $R_{14}[71427]$ & $R_{16}[62499]$ & $R_{16}[62499]$\\
$[\frac{3}{4},1]$ & $Z_2[499995]$ & $R_{10}[99999]$ & $R_{10}[99999]$ & $R_{10}[99999]$ & $R_{10}[99999]$ & $R_{10}[99999]$\\
\bottomrule
\end{tabular}
\end{minipage}
\end{center}
\medskip
\subsection{Finite round ledger}
Within a round all four evaluations read the exponents available at its start.
Each output is rounded upward to twelve decimal places before the next round.
The valid H\"older convex envelope is updated only after the four evaluations.
Every entry below is an exact terminating decimal; $D$ is the strict
record criterion, defined by
\[D=5\Delta_{38}^{(9)}+74\Delta_{40}^{(9)}-11.\]
\begin{center}
\small
\setlength{\tabcolsep}{4pt}
\begin{tabular}{rrrrrr}
\toprule
Round & $\Delta_{34}$ & $\Delta_{36}$ & $\Delta_{38}$ & $\Delta_{40}$ & $D$\\
\midrule
0 & 0.329202900000 & 0.252047100000 & 0.190186000000 & 0.142056900000 & 0.463140600000\\
1 & 0.328601499965 & 0.250127455652 & 0.188302970170 & 0.139862453254 & 0.291336391646\\
2 & 0.328332398257 & 0.249150757590 & 0.186757792931 & 0.138182409445 & 0.159287263585\\
3 & 0.328211986343 & 0.248661104682 & 0.185852461160 & 0.136951673808 & 0.063686167592\\
4 & 0.328158106979 & 0.248505402575 & 0.185380158892 & 0.136178518258 & 0.004111145552\\
5 & 0.328133998186 & 0.248457241158 & 0.185148057878 & 0.135725530085 & -0.030570484320\\
6 & 0.328123210495 & 0.248435690910 & 0.185097318519 & 0.135472843760 & -0.049522969165\\
\bottomrule
\end{tabular}
\end{center}
The final exact criterion is
\[D=-9904593833/200000000000.\]
The convenient upper pair
\[(\Delta_{38},\Delta_{40})\le(1851/10000,271/2000)\]
gives minor-arc saving $19/36000>1/2000$ for 46 variables.
\subsection{Reproduction and file identities}
The parameter JSON and checker form a self-contained two-file replay using only
the Python standard library.  The supplied certificate records every endpoint
evaluation, backward recurrence, exact fraction, and rounded output.
Running the command in an otherwise empty directory reproduces its SHA-256.
\begin{verbatim}
python exact_height_replay.py --fixed \
  --source k9_z2_only_fixed_parameters.json \
  --rounds 6 --out k9_z2_only_exact_certificate.json
\end{verbatim}
SHA-256 identities (the split between lines is typographical):
\begin{quote}\small\ttfamily
\detokenize{k9_z2_only_fixed_parameters.json}\par
7f7cf7c3ae465e11cb5ca273f124f71e\par
736697d5d3bb33c3717723db16e636cb\par\medskip
\detokenize{k9_z2_only_exact_certificate.json}\par
6514323a363a937cfad1951bf26c40cf\par
83b4f54b3df76ad77034a17a3dcb7aa5\par\medskip
\detokenize{exact_height_replay.py}\par
1763b32133010b79f3fe1b09c71f6718\par
b2fa7e21bf33c1e15ee211994a8696c1\par\medskip
\end{quote}
The checker validates the finite arithmetic conditional on the analytic rules proved
in the text.  It does not replace those proofs or the representation argument.

\clearpage

% Source appendix: transfer/appendix_g8_parameters.tex
\section{Complete finite certificate for G(8)}
\label{app:g8-parameters}
This appendix prints the complete fixed recipe corresponding to the file
\par
\path{g8_fixed_parameters.json}.
It uses four policies, 12 finite rounds, and 48 counted updates.
\subsection{Published input exponents}
The following decimals are the exact published upper exponents used as seeds.
The first column of each block is the half-order $s$; the listed value is
$\mu_{2s}^{(d)}$ in $\int_0^1|g_d|^{2s}\ll P^{\mu_{2s}^{(d)}+\varepsilon}$.
All three low-order values and all interpolation operations are independent
of the numerical selection runs: the additional seeds are
$\mu_0^{(d)}=0$, $\mu_2^{(d)}=1$, and $\mu_4^{(d)}=2$.
Only H\"older interpolation between known permissible exponents is used.
\begin{center}
\small
\begin{tabular}{rr@{\hspace{2.2cm}}rr}
\toprule
$s$ & $\mu_{2s}^{(8)}$ & $s$ & $\mu_{2s}^{(16)}$\\
\midrule
3 & 3.0469787 & 3 & 3.0074816\\
4 & 4.2164889 & 4 & 4.0638321\\
5 & 5.4915710 & 5 & 5.1787385\\
6 & 6.8566563 & 6 & 6.3412246\\
7 & 8.3105992 & 7 & 7.5457229\\
8 & 9.8428621 & 8 & 8.7902441\\
9 & 11.4529104 & 9 & 10.0769155\\
10 & 13.1283069 & 10 & 11.4019639\\
11 & 14.8664781 & 11 & 12.7671799\\
12 & 16.6561197 & 12 & 14.1702162\\
13 & 18.4901012 & 13 & 15.6119303\\
14 & 20.3623532 & 14 & 17.0906280\\
15 & 22.2661078 & 15 & 18.6062852\\
16 & 24.1891161 & 16 & 20.1573791\\
17 & 26.1294925 & 17 & 21.7431712\\
18 & 28.0833353 & 18 & 23.3621148\\
19 & 30.0473193 & 19 & 25.0129891\\
20 & 32.0186056 & 20 & 26.6941169\\
 &  & 21 & 28.4039654\\
 &  & 22 & 30.1407641\\
 &  & 23 & 31.9028113\\
 &  & 24 & 33.6883143\\
 &  & 25 & 35.4955330\\
 &  & 26 & 37.3227332\\
 &  & 27 & 39.1682523\\
 &  & 28 & 41.0304795\\
 &  & 29 & 42.9078717\\
 &  & 30 & 44.7989875\\
\bottomrule
\end{tabular}
\end{center}
Source: Vaughan--Wooley, \emph{Further improvements in Waring's problem, IV}, Sections 10 and 18~\cite{VW00}.
The stored pointwise metadata are $\sigma_{8}=0.014356$ and $\sigma_{16}=0.005864$; these two values are not queried by the frozen recipes, whose moment orders use the printed tables or the stated elementary low-order values.
\subsection{The four exact policies}
The path is $(r_1,\ldots,r_j)$ in the counted recurrence.
Each size exponent below has eight decimal places and is exact.
Put $T=\sum_i\phi_i$, $h=k-2-kT$, and $\ell_{\max}=k(1-T)/2$.
A normalized interval $[a,b]$ in the weight tables means
$1+a(\ell_{\max}-1)\le\ell\le1+b(\ell_{\max}-1)$.
\paragraph{Full order $28$.}
$j=6$, $(r_1,\ldots,r_j)=(13,12,11,10,9,8)$.
\[ (\phi_1,\ldots,\phi_j)=\bigl(0.11586958,\ 0.11191889,\ 0.10604341,\ 0.09516681,\ 0.07693878,\ 0.03227716\bigr). \]
$T=0.53821463$, $h=1.69428296$.
The auxiliary terminal interpolates with smooth full order $v=20$,
with $(a,c,\gamma)=(\frac{1}{5},\frac{4}{5},\frac{3}{5})$.
The divisor second-moment prefix conditions are imposed for this policy only:
$T_i+(k-2)\phi_i<1$ for $2\le i\le j$; their least margin is $0.00000045$.
\paragraph{Full order $30$.}
$j=5$, $(r_1,\ldots,r_j)=(14,13,12,11,11)$.
\[ (\phi_1,\ldots,\phi_j)=\bigl(0.11756478,\ 0.11402685,\ 0.10847441,\ 0.09901555,\ 0.08190194\bigr). \]
$T=0.52098353$, $h=1.83213176$.
The terminal uses the direct remaining smooth moment; there is no auxiliary interpolation.
\paragraph{Full order $32$.}
$j=5$, $(r_1,\ldots,r_j)=(15,14,13,12,11)$.
\[ (\phi_1,\ldots,\phi_j)=\bigl(0.11889202,\ 0.11553618,\ 0.11014056,\ 0.10056412,\ 0.08192638\bigr). \]
$T=0.52705926$, $h=1.78352592$.
The terminal uses the direct remaining smooth moment; there is no auxiliary interpolation.
\paragraph{Full order $34$.}
$j=4$, $(r_1,\ldots,r_j)=(16,15,14,14)$.
\[ (\phi_1,\ldots,\phi_j)=\bigl(0.11996738,\ 0.11673903,\ 0.11137032,\ 0.10234991\bigr). \]
$T=0.45042664$, $h=2.39658688$.
The terminal uses the direct remaining smooth moment; there is no auxiliary interpolation.
\subsection{Complete rational H\"older allocations}
An entry $R_p[n]$ means the refined coefficient resource of full order $p$
with weight $n/10^6$.  The entry $B[n]$ is the published basic second-moment
resource with the same weight.  The column $L_i$ identifies the coefficient factor
(so its index is one greater than the zero-based JSON index).
$Z_2[n]$ denotes the sampled quadratic second-moment resource.
$C_p[n]$ denotes the published cube-root coefficient resource.
$Z_4[n]$ denotes the quadratic fourth-moment saving $\min(\ell,2,h)$;
$Z_4^{+}[n]$ denotes its refined concave branch $\min(2\ell-h,2)$.
When two entries occur in one cell, both allocations are used and their
weights are added.  An empty cell has weight zero.  All weights are nonnegative.
The replay checks total H\"older mass at most one and total moment load at most one
for each factor.  These printed integer numerators, together with the preceding
exact size exponents, give the complete policy; no optimizer output is needed.
\begin{center}
\begin{minipage}{\linewidth}\centering
\footnotesize
\setlength{\tabcolsep}{3.3pt}
\renewcommand{\arraystretch}{1.22}
\textbf{Full order $28$: fixed weights}\par\smallskip
\begin{tabular}{cccccccc}
\toprule
$[a,b]$ & $Z$ & $L_1$ & $L_2$ & $L_3$ & $L_4$ & $L_5$ & $L_6$\\
\midrule
$[0,\frac{1}{8}]$ & $Z_2[499995]$ & $R_{44}[22727]$ & $R_{46}[21738]$ & $R_{48}[20833]$ & \shortstack[c]{$B[414750]$\\$R_{52}[3278]$} & $C_{60}[16666]$ & --\\
$[\frac{1}{8},\frac{1}{4}]$ & $Z_2[499995]$ & $R_{40}[24999]$ & $R_{42}[23809]$ & $R_{44}[22727]$ & \shortstack[c]{$B[406538]$\\$R_{46}[4063]$} & $R_{56}[17856]$ & --\\
$[\frac{1}{4},\frac{3}{8}]$ & $Z_2[499995]$ & $R_{38}[26315]$ & $R_{38}[26315]$ & $R_{40}[24999]$ & \shortstack[c]{$B[397482]$\\$R_{42}[4881]$} & $R_{50}[19999]$ & --\\
$[\frac{3}{8},\frac{1}{2}]$ & $Z_2[499995]$ & $R_{34}[29411]$ & $R_{34}[29411]$ & $R_{36}[27777]$ & \shortstack[c]{$B[384593]$\\$R_{38}[6073]$} & $R_{44}[22727]$ & --\\
$[\frac{1}{2},\frac{5}{8}]$ & $Z_2[499995]$ & $R_{30}[33332]$ & $R_{32}[31249]$ & $R_{32}[31249]$ & \shortstack[c]{$B[372875]$\\$R_{34}[7477]$} & $R_{42}[23809]$ & --\\
$[\frac{5}{8},\frac{3}{4}]$ & $Z_2[499995]$ & $R_{28}[35713]$ & $R_{28}[35713]$ & $R_{28}[35713]$ & $R_{32}[31249]$ & \shortstack[c]{$B[354683]$\\$R_{42}[6919]$} & --\\
$[\frac{3}{4},\frac{7}{8}]$ & $Z_2[499995]$ & $R_{24}[41666]$ & $R_{24}[41666]$ & $R_{26}[38461]$ & $R_{32}[31249]$ & \shortstack[c]{$B[339299]$\\$R_{42}[7652]$} & --\\
$[\frac{7}{8},1]$ & \shortstack[c]{$Z_2[135579]$\\$Z_4^{+}[182207]$} & $R_{22}[45454]$ & $R_{24}[41666]$ & $R_{26}[38461]$ & \shortstack[c]{$B[27068]$\\$R_{32}[29557]$} & $B[499995]$ & --\\
\bottomrule
\end{tabular}
\end{minipage}
\end{center}
\medskip
\begin{center}
\begin{minipage}{\linewidth}\centering
\footnotesize
\setlength{\tabcolsep}{3.3pt}
\renewcommand{\arraystretch}{1.22}
\textbf{Full order $30$: fixed weights}\par\smallskip
\begin{tabular}{ccccccc}
\toprule
$[a,b]$ & $Z$ & $L_1$ & $L_2$ & $L_3$ & $L_4$ & $L_5$\\
\midrule
$[0,\frac{1}{8}]$ & $Z_2[499995]$ & $R_{44}[22727]$ & $R_{44}[22727]$ & $R_{46}[21738]$ & \shortstack[c]{$B[412042]$\\$R_{50}[3518]$} & $R_{58}[17241]$\\
$[\frac{1}{8},\frac{1}{4}]$ & $Z_2[499995]$ & $R_{40}[24999]$ & $R_{40}[24999]$ & $R_{42}[23809]$ & \shortstack[c]{$B[402525]$\\$R_{44}[4430]$} & $R_{52}[19230]$\\
$[\frac{1}{4},\frac{3}{8}]$ & $Z_2[499995]$ & $R_{36}[27777]$ & $R_{38}[26315]$ & $R_{38}[26315]$ & \shortstack[c]{$B[392471]$\\$R_{40}[5376]$} & $R_{46}[21738]$\\
$[\frac{3}{8},\frac{1}{2}]$ & $Z_2[499995]$ & $R_{32}[31249]$ & $R_{34}[29411]$ & $R_{34}[29411]$ & \shortstack[c]{$B[378153]$\\$R_{36}[6768]$} & $R_{40}[24999]$\\
$[\frac{1}{2},\frac{5}{8}]$ & $Z_2[499995]$ & $R_{28}[35713]$ & $R_{30}[33332]$ & $R_{30}[33332]$ & \shortstack[c]{$B[361160]$\\$R_{32}[8677]$} & $R_{36}[27777]$\\
$[\frac{5}{8},\frac{3}{4}]$ & $Z_2[499995]$ & $R_{24}[41666]$ & $R_{26}[38461]$ & $R_{26}[38461]$ & $R_{28}[35713]$ & \shortstack[c]{$B[336615]$\\$R_{36}[9076]$}\\
$[\frac{3}{4},\frac{7}{8}]$ & $Z_2[499995]$ & $R_{20}[49999]$ & $R_{20}[49999]$ & $R_{22}[45454]$ & $R_{26}[38461]$ & \shortstack[c]{$B[305262]$\\$R_{36}[10818]$}\\
$[\frac{7}{8},1]$ & \shortstack[c]{$Z_2[105125]$\\$Z_4^{+}[197434]$} & $R_{18}[55554]$ & $R_{20}[49999]$ & $R_{22}[45454]$ & \shortstack[c]{$B[8628]$\\$R_{26}[37797]$} & $B[499994]$\\
\bottomrule
\end{tabular}
\end{minipage}
\end{center}
\medskip
\begin{center}
\begin{minipage}{\linewidth}\centering
\footnotesize
\setlength{\tabcolsep}{3.3pt}
\renewcommand{\arraystretch}{1.22}
\textbf{Full order $32$: fixed weights}\par\smallskip
\begin{tabular}{ccccccc}
\toprule
$[a,b]$ & $Z$ & $L_1$ & $L_2$ & $L_3$ & $L_4$ & $L_5$\\
\midrule
$[0,\frac{1}{8}]$ & $Z_2[499995]$ & $R_{44}[22727]$ & $R_{44}[22727]$ & $R_{46}[21738]$ & \shortstack[c]{$B[411889]$\\$R_{48}[3671]$} & $R_{58}[17241]$\\
$[\frac{1}{8},\frac{1}{4}]$ & $Z_2[499995]$ & $R_{40}[24999]$ & $R_{40}[24999]$ & $R_{42}[23809]$ & \shortstack[c]{$B[402525]$\\$R_{44}[4430]$} & $R_{52}[19230]$\\
$[\frac{1}{4},\frac{3}{8}]$ & $Z_2[499995]$ & $R_{36}[27777]$ & $R_{38}[26315]$ & $R_{38}[26315]$ & \shortstack[c]{$B[392471]$\\$R_{40}[5376]$} & $R_{46}[21738]$\\
$[\frac{3}{8},\frac{1}{2}]$ & $Z_2[499995]$ & $R_{32}[31249]$ & $R_{34}[29411]$ & $R_{34}[29411]$ & \shortstack[c]{$B[379414]$\\$R_{36}[6698]$} & $R_{42}[23809]$\\
$[\frac{1}{2},\frac{5}{8}]$ & $Z_2[499995]$ & $R_{30}[33332]$ & $R_{30}[33332]$ & $R_{30}[33332]$ & \shortstack[c]{$B[363700]$\\$R_{32}[8518]$} & $R_{36}[27777]$\\
$[\frac{5}{8},\frac{3}{4}]$ & $Z_2[499995]$ & $R_{26}[38461]$ & $R_{26}[38461]$ & $R_{26}[38461]$ & $R_{28}[35713]$ & \shortstack[c]{$B[340009]$\\$R_{36}[8888]$}\\
$[\frac{3}{4},\frac{7}{8}]$ & $Z_2[499995]$ & $R_{22}[45454]$ & $R_{22}[45454]$ & $R_{22}[45454]$ & $R_{26}[38461]$ & \shortstack[c]{$B[314887]$\\$R_{36}[10283]$}\\
$[\frac{7}{8},1]$ & \shortstack[c]{$Z_2[106167]$\\$Z_4^{+}[196913]$} & $R_{20}[49999]$ & $R_{20}[49999]$ & $R_{22}[45454]$ & \shortstack[c]{$B[14082]$\\$R_{26}[37377]$} & $B[499994]$\\
\bottomrule
\end{tabular}
\end{minipage}
\end{center}
\medskip
\begin{center}
\begin{minipage}{\linewidth}\centering
\footnotesize
\setlength{\tabcolsep}{3.3pt}
\renewcommand{\arraystretch}{1.22}
\textbf{Full order $34$: fixed weights}\par\smallskip
\begin{tabular}{cccccc}
\toprule
$[a,b]$ & $Z$ & $L_1$ & $L_2$ & $L_3$ & $L_4$\\
\midrule
$[0,\frac{1}{8}]$ & $Z_2[499995]$ & \shortstack[c]{$B[428305]$\\$R_{42}[3413]$} & $R_{42}[23809]$ & $R_{44}[22727]$ & $R_{46}[21738]$\\
$[\frac{1}{8},\frac{1}{4}]$ & $Z_2[499995]$ & \shortstack[c]{$B[420696]$\\$R_{38}[4173]$} & $R_{38}[26315]$ & $R_{40}[24999]$ & $R_{42}[23809]$\\
$[\frac{1}{4},\frac{3}{8}]$ & $Z_2[499995]$ & \shortstack[c]{$B[407621]$\\$R_{32}[5773]$} & $R_{34}[29411]$ & $R_{34}[29411]$ & $R_{36}[27777]$\\
$[\frac{3}{8},\frac{1}{2}]$ & $Z_2[499995]$ & \shortstack[c]{$B[389739]$\\$R_{28}[7875]$} & $R_{28}[35713]$ & $R_{30}[33332]$ & $R_{30}[33332]$\\
$[\frac{1}{2},\frac{5}{8}]$ & $Z_2[499995]$ & \shortstack[c]{$B[358329]$\\$R_{22}[12878]$} & $R_{22}[45454]$ & $R_{24}[41666]$ & $R_{24}[41666]$\\
$[\frac{5}{8},\frac{3}{4}]$ & $Z_2[499995]$ & \shortstack[c]{$B[260414]$\\$R_{14}[34225]$} & $R_{14}[71427]$ & $R_{14}[71427]$ & $R_{16}[62499]$\\
$[\frac{3}{4},\frac{7}{8}]$ & \shortstack[c]{$Z_2[466782]$\\$Z_4[16606]$} & \shortstack[c]{$R_{6}[66425]$\\$R_{8}[75179]$} & $R_{8}[124998]$ & $R_{8}[124998]$ & $R_{8}[124998]$\\
$[\frac{7}{8},1]$ & \shortstack[c]{$Z_4[69863]$\\$Z_4^{+}[180133]$} & $R_{4}[249997]$ & $R_{6}[166664]$ & $R_{6}[166664]$ & $R_{6}[166664]$\\
\bottomrule
\end{tabular}
\end{minipage}
\end{center}
\medskip
\subsection{Finite round ledger}
Within a round all four evaluations read the exponents available at its start.
Each output is rounded upward to twelve decimal places before the next round.
The valid H\"older convex envelope is updated only after the four evaluations.
Every entry below is an exact terminating decimal; $D$ is the strict
record criterion, defined by
\[D=\Delta_{32}^{(8)}+7\Delta_{34}^{(8)}-1.\]
\begin{center}
\small
\setlength{\tabcolsep}{4pt}
\begin{tabular}{rrrrrr}
\toprule
Round & $\Delta_{28}$ & $\Delta_{30}$ & $\Delta_{32}$ & $\Delta_{34}$ & $D$\\
\midrule
0 & 0.362353200000 & 0.266107800000 & 0.189116100000 & 0.129492500000 & 0.095563600000\\
1 & 0.361191319429 & 0.262218424695 & 0.185605866121 & 0.126167376968 & 0.068777504897\\
2 & 0.360677692450 & 0.260279164129 & 0.182620068594 & 0.123078256650 & 0.044167865144\\
3 & 0.360450635832 & 0.259324859400 & 0.180891421884 & 0.120528807443 & 0.024593073985\\
4 & 0.360350262001 & 0.258860185296 & 0.179928609981 & 0.119040136099 & 0.013209562674\\
5 & 0.360305890222 & 0.258635880876 & 0.179407280490 & 0.118185010380 & 0.006702353150\\
6 & 0.360286275002 & 0.258528389904 & 0.179131041086 & 0.117702043622 & 0.003045346440\\
7 & 0.360277603796 & 0.258477194870 & 0.178987162647 & 0.117433871477 & 0.001024262986\\
8 & 0.360273770557 & 0.258452940962 & 0.178913267636 & 0.117287113684 & -0.000076936576\\
9 & 0.360272076016 & 0.258441503338 & 0.178875756814 & 0.117207861261 & -0.000669214359\\
10 & 0.360271326918 & 0.258436131332 & 0.178856903359 & 0.117165576462 & -0.000984061407\\
11 & 0.360270995768 & 0.258434347629 & 0.178847626533 & 0.117143260275 & -0.001149551542\\
12 & 0.360270849378 & 0.258434055411 & 0.178844726003 & 0.117131705619 & -0.001233334664\\
\bottomrule
\end{tabular}
\end{center}
The final exact criterion is
\[D=-154166833/125000000000.\]
The convenient upper pair
\[(\Delta_{32},\Delta_{34})\le(1789/10000,11721/100000)\]
gives minor-arc saving $7/100000$ for 38 variables.
\subsection{Reproduction and file identities}
The parameter JSON and checker form a self-contained two-file replay using only
the Python standard library.  The supplied certificate records every endpoint
evaluation, backward recurrence, exact fraction, and rounded output.
Running the command in an otherwise empty directory reproduces its SHA-256.
\begin{verbatim}
python g8_standalone_replay.py \
  --source g8_fixed_parameters.json \
  --rounds 12 --out g8_exact_certificate.json
\end{verbatim}
SHA-256 identities (the split between lines is typographical):
\begin{quote}\small\ttfamily
\detokenize{g8_fixed_parameters.json}\par
7af6e60058a3522ab38b71161d365f17\par
73ce5fd4c36cb654532ae21b858de583\par\medskip
\detokenize{g8_exact_certificate.json}\par
be95505ae57e69dc915cc83e7aa99aba\par
7e91d9d89aa6e63083355f6582ac7983\par\medskip
\detokenize{g8_standalone_replay.py}\par
62844204b0d06c1ac142b9fd5cfc1378\par
ad9cb9f33055b3bf541f9dccf621bcc3\par\medskip
\end{quote}
The checker validates the finite arithmetic conditional on the analytic rules proved
in the text.  It does not replace those proofs or the representation argument.

\clearpage

\clearpage
\begin{thebibliography}{20}\small
\bibitem{VW94} R. C. Vaughan and T. D. Wooley,
\emph{Further improvements in Waring's problem, II: sixth powers},
Duke Math. J. 76 (1994), 683--710.
\url{https://doi.org/10.1215/S0012-7094-94-07626-6}.
\bibitem{VW95} R. C. Vaughan and T. D. Wooley,
\emph{Further improvements in Waring's problem},
Acta Math. 174 (1995), 147--240.
\url{https://www.math.purdue.edu/~twooley/publ/1995%20fiwp.pdf}.
\bibitem{VW00} R. C. Vaughan and T. D. Wooley,
\emph{Further improvements in Waring's problem, IV: higher powers},
Acta Arith. 94 (2000), 203--285.
\url{https://www.math.purdue.edu/~twooley/publ/2000%20fiwp4.pdf}.
\bibitem{W93} T. D. Wooley,
\emph{The application of a new mean value theorem to the fractional parts
of polynomials}, Acta Arith. 65 (1993), 163--179.
\url{https://doi.org/10.4064/aa-65-2-163-179}.
\bibitem{W00QD} T. D. Wooley,
\emph{Quasi-diagonal behaviour and smooth Weyl sums},
Monatsh. Math. 130 (2000), 161--170.
\url{https://www.math.purdue.edu/~twooley/publ/2000SWSQD.pdf}.
\bibitem{W16} T. D. Wooley,
\emph{On Waring's problem for intermediate powers},
Acta Arith. 176 (2016), 241--247.
\url{https://arxiv.org/abs/1602.03221}.
\bibitem{BW22PN} J. Br\"udern and T. D. Wooley,
\emph{Partitio Numerorum: sums of a prime and a number of $k$-th powers},
2022, arXiv:2211.10387. \url{https://arxiv.org/abs/2211.10387}.
\bibitem{BW23} J. Br\"udern and T. D. Wooley,
\emph{On Waring's problem for larger powers},
J. reine angew. Math. 805 (2023), 115--142.
\url{https://www.math.purdue.edu/~twooley/publ/20230915wpk.pdf}.
\bibitem{BW24} J. Br\"udern and T. D. Wooley,
\emph{Estimates for smooth Weyl sums on major arcs},
Internat. Math. Res. Notices 2024 (2024), no. 24, 14662--14688.
\url{https://doi.org/10.1093/imrn/rnae252}.
\bibitem{Bou16} J. Bourgain, \emph{On the Vinogradov mean value}, 2016,
Theorem 10 and Section 6. \url{https://arxiv.org/abs/1601.08173}.
\bibitem{DKW05} J.-M. Deshouillers, K. Kawada and T. D. Wooley,
\emph{On sums of sixteen biquadrates},
M\'em. Soc. Math. France, nouvelle s\'erie, 100 (2005), vi+120 pp.
\url{https://www.numdam.org/item/MSMF_2005_2_100__1_0/}.
\end{thebibliography}
\end{document}
